\documentclass[11pt]{article}
\usepackage[margin=1in]{geometry}
\usepackage{amsmath,amssymb,amsthm,mathtools}
\IfFileExists{lmodern.sty}{\usepackage{lmodern}}{}
\usepackage{microtype}
\usepackage{booktabs}
\usepackage{enumitem}
\usepackage{xcolor}
\usepackage{tikz-cd}
\usepackage[colorlinks=true,linkcolor=blue!50!black,citecolor=green!35!black,urlcolor=blue!50!black]{hyperref}
\usepackage[nameinlink]{cleveref}

\newtheorem{theorem}{Theorem}[section]
\newtheorem{proposition}[theorem]{Proposition}
\newtheorem{lemma}[theorem]{Lemma}
\newtheorem{corollary}[theorem]{Corollary}
\theoremstyle{definition}
\newtheorem{definition}[theorem]{Definition}
\newtheorem{axiom}[theorem]{Axiom}
\newtheorem{example}[theorem]{Example}
\theoremstyle{remark}
\newtheorem{remark}[theorem]{Remark}

\newcommand{\D}{\mathsf{D}}
\newcommand{\Set}{\mathbf{Set}}
\newcommand{\Hyp}{\mathrm{Hyp}}
\newcommand{\Sh}{\mathrm{Sh}}
\newcommand{\Post}{\mathrm{Post}}
\newcommand{\Cond}{\mathrm{Cond}}
\newcommand{\C}{\mathsf{C}}
\newcommand{\E}{\mathsf{E}}
\newcommand{\id}{\mathrm{id}}
\newcommand{\hid}{h_{\mathrm{id}}}
\newcommand{\lb}{\bar\lambda}
\newcommand{\ls}{\lambda^{\sigma}}
\newcommand{\N}{\mathcal{N}}
\newcommand{\cp}{\circledast}

\title{The Three-Level Posterior Tower:\\
Hierarchical Bayesian Inference and Mixed Distributive Laws\\[4pt]
\large Revised version 5}
\author{Vinay K. Awasthi\\
Department of Applied Mathematics and Statistics\\
Johns Hopkins University\thanks{Code and verification scripts:
\url{https://github.com/vkawasth/Trainingless-Transformer-NoGradientDescent/tree/main/Heirarchical_posteriors}}}
\date{October 3, 2026}

\begin{document}
\maketitle

\begin{abstract}
We construct a categorical framework for structured Bayesian inference that keeps three levels of
posterior information as state: a posterior $\rho$ over outcomes, a second-order posterior $\Pi$ over
outcome posteriors, a third-order posterior $\Omega$, and a shadow $\Sigma$ over structured hypotheses.
Combining this state with the finite distribution monad $\D$ forces a design choice. In the
\emph{absolute} tower $\C_3(X)=X\times\D X\times\D^2X\times\D^3X\times\Sh(X)$ the levels are
re-indexed by the object, so $\C_3(\D X)$ already contains $\D^4X$: the hierarchy grows by one order
every time the monad is applied. We prove that $X\times\E(X)$ is a comonad for every endofunctor $\E$,
that the context and distribution-level lifts are coherent by functoriality, and that the two natural
candidate laws $\D\C_3\Rightarrow\C_3\D$ each fail at least one Beck axiom; the unit axiom alone forbids
point-mass hierarchy at level one. We then construct a \emph{base-relative} tower, in which the three
levels live over a fixed outcome base $B$ and do not grow. Its \emph{coherent} contexts, those whose
barycentric projections and shadow predictive agree, form a sub-algebra of the monad, and the
canonical hierarchy $\rho_h$, $\Pi_h=\delta_{\rho_h}$, $\Omega_h=\delta_{\Pi_h}$ extends to an affine
section of it. For this tower we obtain a strict mixed distributive law satisfying all four Beck
axioms, natural in the base. Pooling commutes with hierarchy formation, conditioning commutes with
pooling after evidence reweighting, and distinct posteriors are retained at level two. We then pose a coherent lifting problem: given
admissible first- and second-order states generated by a finite hypothesis class, when does a claimed
posterior state extend to a full coherent three-level context? Admissibility at each level is a
convex condition, failure is monotone up the tower, each level can be the first to fail, and every
failure carries an explicit certificate: a separating test function at level one (with the obstruction
size given by a minimax identity) and a witness state at levels two and three. Pooling can hide
level-one obstructions but never level-two or level-three ones. This gives a precise candidate
definition of a hallucinated claim as one whose required coherent lift is obstructed. We identify
the coalgebras of the base-relative comonad with sets over the context algebra and its bialgebras
with convex spaces over it, and we make precise that the ambient 2-category is $\mathbf{Cat}$, in which
these categories are locally discrete, so two-dimensional monad theory adds no content here; at level
one the defects of the absolute shift law are ordered in the convex order, which locates where genuine
2-cells could arise. On the
algebraic side we give a normalized partial composition on structured posteriors, corrected
conditioning laws, an exact characterization of when observational equivalence is a congruence, and
a proof that the likelihood-factorization ``iff'' holds under nondegeneracy alone. All identities and
counterexamples were checked in exact rational arithmetic; the scripts are in the accompanying
repository.
\end{abstract}

\tableofcontents

%=====================================================================
\section{Introduction}
%=====================================================================

Bayesian inference in structured problems produces a hierarchy of posteriors: a posterior over
outcomes, a posterior over hypotheses (each of which itself predicts outcomes), and a posterior over
the ways hypotheses combine. This paper builds a categorical framework that retains all three levels
as state rather than as intermediate computations, and asks how such state interacts with the
probability monad.

The framework builds on the three-layer architecture of~\cite{Awasthi-ThreeLayers}; the hierarchy
layer is developed in~\cite{Awasthi-Hierarchy} and the holonomy obstruction
in~\cite{Awasthi-Holonomy}.

\paragraph{The design choice.} The levels of the tower are indexed by the object on which the
comonad is evaluated. Evaluating at $\D X$, as any distributive law $\D\C_3\Rightarrow\C_3\D$ must,
shifts every level up by one order:
\[
\C_3(\D X)=\D X\times\D^2X\times\D^3X\times\D^4X\times\Sh(\D X).
\]
There are two ways to proceed. \emph{Option~1} keeps this absolute indexing and accepts a target that
contains a fourth-order level (\cref{sec:laws}); the transformation of the first version omitted that
level and was therefore ill-typed. \emph{Option~2} fixes an outcome base $B$ once and for all and lets
the three levels live over $B$ whatever object the comonad is evaluated at (\cref{sec:relative}).
Option~2 preserves exactly three hierarchical levels and is the main construction of this paper; it
admits a strict mixed distributive law (\cref{thm:relative-law}). Option~1 is developed to show
precisely what goes wrong without a fixed base.

\paragraph{Contributions.}
\begin{enumerate}[leftmargin=*]
\item A base-relative three-level hierarchy over a fixed base $B$, with coherent contexts, the
  canonical hierarchy $\sigma(\delta_h)=(\rho_h,\delta_{\rho_h},\delta_{\delta_{\rho_h}},\delta_h)$ as an
  affine section, and a strict mixed distributive law satisfying all four Beck axioms, natural in the
  base (\cref{sec:relative}).
\item A coherent lifting problem for claims against the tower, with level-wise obstruction
  certificates, strictness and monotonicity of failure levels, and the behavior of obstructions under
  pooling (\cref{sec:lifting}).
\item The coalgebras, lifted monad and bialgebras of the base-relative law, a precise statement of
  the ambient 2-category and of what two-dimensional monad theory does and does not add, and a
  convex-order comparison of the shift law's defects (\cref{sec:coalg}).
\item A general \emph{context comonad} $\C_\E(X)=X\times\E(X)$ for any endofunctor $\E$ of $\Set$
  (\cref{thm:context-comonad}), specialized to the three-level tower $\C_3$.
\item A hypothesis functor $(\Hyp,q)$ with $q$ natural, so that the compatibility square is
  naturality (\cref{sec:hyp}); the context lift $J=\Hyp(j)$ and the distribution-level lift $K$
  satisfy all coherence conditions by proof (\cref{prop:CL,prop:K}).
\item A normalized partial algebra $(\oplus,\otimes,\Cond)$ on structured posteriors with corrected
  laws (\cref{sec:post}).
\item An affine first isomorphism theorem for observational equivalence, and an exact
  characterization of when it is a congruence for $\otimes$ and for $\Cond$ (\cref{sec:obs}).
\item For the absolute tower, two mixed distributive laws $\ls$ and $\lb$ with a complete account
  of which Beck axioms each satisfies (\cref{sec:laws}, summarized in \cref{tab:status}), together with a structural lemma
  showing that the unit axiom forbids point-mass hierarchy at level one (\cref{lem:unit-forces}).
\item A proof that composite identifiability trivializes the hypothesis space, and a proof that the
  likelihood-factorization ``iff'' holds under nondegeneracy alone (\cref{rem:ci,thm:iff}).
\end{enumerate}

\paragraph{Relation to earlier versions.} The first version claimed that a three-level
transformation $\lambda^{[3]}$ satisfies the unit, counit and comultiplication axioms exactly and
extends to a weak distributive law with a barycentric 2-cell. That transformation omitted the
$\D^4X$ component of its target, and the claims do not hold; the corrected statements for the
absolute tower are \cref{thm:shift,thm:avg,thm:avg-comult}. The second version corrected the typing
but left the existence of a full law open. Version 3 resolved the question for the base-relative
tower; version 4 added the lifting problem; version 5 corrects the 2-categorical discussion and
several empirical attributions. A full list of changes is in \cref{app:changes}.

%=====================================================================
\section{Preliminaries}
%=====================================================================

\subsection{The finite distribution monad}

Let $\D$ be the finitely supported distribution monad on $\Set$: $\D(X)$ is the set of finitely
supported probability measures on $X$, $\D(f)=f_*$ is pushforward, $\eta_X(x)=\delta_x$, and
$\mu_X\colon \D^2X\to\D X$ is barycenter, $\mu_X(\sum_i a_i\delta_{\rho_i})=\sum_i a_i\rho_i$. Its
Kleisli category restricted to finite sets is the category of finite stochastic matrices; see
\cite{Fritz2020,FGPR2023} for the Markov-category perspective.

We work on $\Set$ rather than on finite sets: even for finite $X$ with $|X|\ge 2$, the set $\D(X)$ is
an infinite simplex, so $\D$ does not restrict to an endofunctor of finite sets. Finiteness enters
only through finite supports.

Each $\D(Y)$ is a convex set, and every $\D(f)$ is affine. A product of convex sets is convex
componentwise, and a product of affine maps is affine.

\subsection{Mixed distributive laws}

\begin{definition}[Mixed distributive law~\cite{Beck1969,PowerWatanabe2002}]\label{def:mixed}
Let $(T,\eta,\mu)$ be a monad and $(C,\varepsilon,\delta)$ a comonad on the same category. A
\emph{mixed distributive law} is a natural transformation $\lambda\colon TC\Rightarrow CT$ such that
\begin{align*}
\text{(U)}\quad & \lambda\circ\eta_{C} = C\eta, &
\text{(M)}\quad & \lambda\circ\mu_{C} = C\mu\circ\lambda_{T}\circ T\lambda,\\
\text{(Co)}\quad & \varepsilon_{T}\circ\lambda = T\varepsilon, &
\text{(Cm)}\quad & \delta_{T}\circ\lambda = C\lambda\circ\lambda_{C}\circ T\delta.
\end{align*}
\end{definition}

Weak distributive laws between two monads, in the sense of Street and B\"ohm
\cite{Street2009,Bohm2010,BohmLackStreet2011}, keep the multiplication-type axioms and relax the
unit-type ones; Garner~\cite{Garner2020} and Goy--Petri\c{s}an~\cite{GoyPetrisan2020} give
probabilistic and topological examples. We return to this comparison in \cref{rem:weak}.

%=====================================================================
\section{Structured hypotheses}\label{sec:hyp}
%=====================================================================

\begin{definition}[Hypothesis functor]\label{def:hypfunctor}
A \emph{hypothesis functor} is a functor $\Hyp\colon\Set\to\Set$ with a natural transformation
$q\colon\Hyp\Rightarrow\D$, the \emph{predictive map}. Naturality of $q$ is the compatibility square
\[
\begin{tikzcd}
\Hyp(X) \arrow[r,"q_X"] \arrow[d,"\Hyp(f)"'] & \D(X) \arrow[d,"f_*"]\\
\Hyp(Y) \arrow[r,"q_Y"'] & \D(Y)
\end{tikzcd}
\qquad\text{i.e.}\qquad q_Y\circ\Hyp(f)=f_*\circ q_X .
\]
\end{definition}

\begin{definition}[Structured hypotheses]\label{def:structured}
Fix a set $\Theta$ of \emph{structural descriptors}
\[\theta=(I_\theta,F_\theta,E_\theta,L_\theta,G_\theta,C_\theta):\]
a sparse multi-index of
directions, a factor graph, an evidence set, spline links, geometric data and conditional structure.
Put
\[
\Hyp(X)=\Theta\times\D(X),\qquad \Hyp(f)(\theta,\rho)=(\theta,f_*\rho),\qquad q_X(\theta,\rho)=\rho .
\]
\end{definition}

\begin{lemma}\label{lem:hyp-functor}
$\Hyp$ is a functor and $q$ is natural.
\end{lemma}
\begin{proof}
$\Hyp=\Theta\times\D(-)$ is a product of a constant functor with $\D$, and $q$ is the projection,
which is natural because $q_Y(\Hyp(f)(\theta,\rho))=f_*\rho=f_*(q_X(\theta,\rho))$.
\end{proof}

\begin{remark}[Why descriptors are not transported]\label{rem:variance}
Directions $v_j$ and spline links $s_j$ are functions \emph{on} $X$. Such data pull back along maps
\emph{into} $X$; they cannot be pushed along an arbitrary $f\colon X\to Y$ (the expression $v_j\circ f$
is not defined), and conditional structure has no canonical pushforward along non-injective maps. In
CausalIDView the descriptors parametrize a model on a fixed design, and the only datum that must move
along outcome maps is the predictive distribution. \Cref{def:structured} records exactly this. The
\emph{admissible} hypotheses $\Hyp^{\mathrm{adm}}(X)=\{(\theta,\rho_\theta)\}$, in which the
predictive is the one produced by $\theta$, form a subset that is generally not closed under
$\Hyp(f)$; all constructions below are carried out on the full functor $\Hyp$.
\end{remark}

\begin{definition}[Shadow]
$\Sh(X)=\D(\Hyp(X))$, the finite convex combinations $\Sigma=\sum_i\lambda_i\delta_{h_i}$ of
hypotheses. $\Sh$ is a functor with $\Sh(f)=\D(\Hyp(f))$.
\end{definition}

%=====================================================================
\section{The structured posterior}\label{sec:post}
%=====================================================================

\begin{definition}
$\Post(X)=\D(X\times\Hyp(X))$. An element is $P=\sum_k\mu_k\delta_{(x_k,h_k)}$. The \emph{outcome
marginal} is $\pi_{X*}=\D(\mathrm{pr}_1)\colon\Post(X)\to\D(X)$, and the \emph{predictive} is
$p_X=\mu_X\circ\D(q_X\circ\mathrm{pr}_2)$, so $p_X(P)=\sum_k\mu_k\,q_X(h_k)$.
\end{definition}

For a nonzero finitely supported measure $m\ge 0$ write $\N(m)=m/\lVert m\rVert$ for its
normalization. $\N$ absorbs positive scalars: $\N(cm)=\N(m)$ for $c>0$.

\subsection{Mixture}

\begin{definition}
$P_1\oplus_\lambda P_2=\lambda P_1+(1-\lambda)P_2$ for $\lambda\in[0,1]$.
\end{definition}

\begin{lemma}[Barycentric associativity]\label{lem:baryassoc}
For $\lambda,\mu\in[0,1]$ with $\lambda\mu<1$ and $\nu=\frac{(1-\lambda)\mu}{1-\lambda\mu}$,
\[
(P_1\oplus_\lambda P_2)\oplus_\mu P_3 = P_1\oplus_{\lambda\mu}(P_2\oplus_\nu P_3).
\]
If $\lambda=\mu=1$ both sides equal $P_1$ for any choice of the inner weight.
\end{lemma}
\begin{proof}
The left side is $\lambda\mu P_1+(1-\lambda)\mu P_2+(1-\mu)P_3$. On the right,
$(1-\lambda\mu)\nu=(1-\lambda)\mu$ and $(1-\lambda\mu)(1-\nu)=1-\lambda\mu-(1-\lambda)\mu=1-\mu$. The
boundary case is immediate.
\end{proof}

\subsection{Composition}

\begin{definition}[Compositional structure]\label{def:comp}
A compositional structure on $\Hyp(X)$ consists of a relation $\perp$ (compatibility), a map
$\otimes$ defined on compatible pairs, and an element $\hid$, such that:
\begin{enumerate}[label=(A\arabic*)]
\item (ternary admissibility and associativity) $h_1\perp h_2$ and $(h_1\otimes h_2)\perp h_3$ if and
  only if $h_2\perp h_3$ and $h_1\perp(h_2\otimes h_3)$, and then
  $(h_1\otimes h_2)\otimes h_3=h_1\otimes(h_2\otimes h_3)$;
\item (identity) $\hid\perp h$, $h\perp\hid$ and $\hid\otimes h=h\otimes\hid=h$ for all $h$.
\end{enumerate}
\end{definition}

\begin{definition}[Composition of posteriors]
The \emph{unnormalized product} is
\[
(P_1\cp P_2)(x,h)=\sum_{\substack{h_1\perp h_2\\ h_1\otimes h_2=h}}P_1(x,h_1)\,P_2(x,h_2),
\]
a finitely supported measure of total mass $Z(P_1,P_2)\in[0,1]$. If $Z(P_1,P_2)>0$ put
$P_1\otimes P_2=\N(P_1\cp P_2)\in\Post(X)$; otherwise $P_1\otimes P_2$ is undefined.
\end{definition}

\begin{remark}
Without normalization the product is not a probability distribution: if $\perp$ is total its mass
is $\sum_x m_1(x)m_2(x)$, where $m_i=\pi_{X*}(P_i)$, which is strictly less than $1$ unless the
marginals are equal point masses.
\end{remark}

\begin{lemma}\label{lem:assoc}
$\cp$ extends bilinearly to finitely supported measures and is associative. Consequently
$(P_1\otimes P_2)\otimes P_3=P_1\otimes(P_2\otimes P_3)$, and either side is defined exactly when
$P_1\cp P_2\cp P_3\neq 0$.
\end{lemma}
\begin{proof}
Bilinearity is clear from the formula. Expanding both triple products gives sums over triples with
$(h_1\otimes h_2)\otimes h_3=h$, respectively $h_1\otimes(h_2\otimes h_3)=h$, over the admissible
triples; by (A1) these index sets and terms coincide. Since $\cp$ is bilinear and $\N$ absorbs
positive scalars, $\N(\N(m)\cp m')=\N(m\cp m')$, so both normalized products equal
$\N(P_1\cp P_2\cp P_3)$. If $P_1\cp P_2=0$ the triple product is $0$, which gives the definedness
statement.
\end{proof}

\begin{proposition}[Distributivity]\label{prop:distrib}
Let $Z_{i3}=Z(P_i,P_3)$ and suppose $Z_{13},Z_{23}>0$. Then
\[
(P_1\oplus_\lambda P_2)\otimes P_3=(P_1\otimes P_3)\oplus_{\lambda''}(P_2\otimes P_3),\qquad
\lambda''=\frac{\lambda Z_{13}}{\lambda Z_{13}+(1-\lambda)Z_{23}} .
\]
\end{proposition}
\begin{proof}
By bilinearity $(\lambda P_1+(1-\lambda)P_2)\cp P_3=\lambda Z_{13}(P_1\otimes P_3)+(1-\lambda)Z_{23}
(P_2\otimes P_3)$; normalize.
\end{proof}

The weight is $\lambda''=\lambda$ only when $Z_{13}=Z_{23}$ (or $\lambda\in\{0,1\}$).

\subsection{Conditioning}

\begin{definition}
For evidence $e$ with likelihood $L(e\mid\cdot)\colon\Hyp(X)\to[0,\infty)$, put
$Z_e(P)=\sum_{(x,h)}L(e\mid h)P(x,h)$ and, when $Z_e(P)>0$,
$\Cond_e(P)=\N\big(L(e\mid\cdot)\,P\big)$.
\end{definition}

\begin{proposition}[Chain rule]
Suppose $L(e_1\wedge e_2\mid h)=L(e_1\mid h)\,L(e_2\mid h,e_1)$, and write $\Cond_{e_2\mid e_1}$ for
conditioning with likelihood $L(e_2\mid\cdot,e_1)$. Then
$\Cond_{e_2\mid e_1}(\Cond_{e_1}P)=\Cond_{e_1\wedge e_2}(P)$ whenever the right side is defined. In
particular, if $e_1,e_2$ are conditionally independent given $h$, then
$\Cond_{e_2}\circ\Cond_{e_1}=\Cond_{e_1\wedge e_2}$.
\end{proposition}
\begin{proof}
$\N(L_2\cdot\N(L_1P))=\N(L_2L_1P)$.
\end{proof}

\begin{proposition}[Conditioning and mixture]
If $Z_1=Z_e(P_1)>0$ and $Z_2=Z_e(P_2)>0$, then
$\Cond_e(P_1\oplus_\lambda P_2)=\Cond_e(P_1)\oplus_{\lambda'}\Cond_e(P_2)$ with
$\lambda'=\lambda Z_1/(\lambda Z_1+(1-\lambda)Z_2)$.
\end{proposition}
\begin{proof}
$L\cdot(\lambda P_1+(1-\lambda)P_2)=\lambda Z_1\Cond_e(P_1)+(1-\lambda)Z_2\Cond_e(P_2)$; normalize.
\end{proof}

\begin{proposition}[Conditioning and composition]\label{prop:cond-comp}
If $L(e\mid h_1\otimes h_2)=L(e\mid h_1)L(e\mid h_2)$ for all compatible pairs, then
$\Cond_e(P_1\otimes P_2)=\Cond_e(P_1)\otimes\Cond_e(P_2)$, and the left side is defined exactly when
the right side is.
\end{proposition}
\begin{proof}
Pointwise, $L(e\mid h_1\otimes h_2)P_1(x,h_1)P_2(x,h_2)=(LP_1)(x,h_1)(LP_2)(x,h_2)$, so
$L\cdot(P_1\cp P_2)=(LP_1)\cp(LP_2)$. Using that $\N$ absorbs positive scalars and $\cp$ is
bilinear, $\Cond_e(P_1\otimes P_2)=\N(L\cdot(P_1\cp P_2))=\N((LP_1)\cp(LP_2))
=\N(\Cond_eP_1\cp\Cond_eP_2)$. Both sides are defined iff $(LP_1)\cp(LP_2)\ne0$.
\end{proof}

\begin{remark}[Partial algebra]
$\Post(X)$ is a convex set with a partial, normalized, associative composition and a partial
conditioning action. It is not a semiring: $\oplus$ distributes over $\otimes$ only with the
reweighting of \cref{prop:distrib}.
\end{remark}

%=====================================================================
\section{Observational equivalence}\label{sec:obs}
%=====================================================================

\begin{definition}
$P_1\equiv P_2$ iff $\pi_{X*}(P_1)=\pi_{X*}(P_2)$.
\end{definition}

\begin{theorem}[Affine first isomorphism theorem]\label{thm:firstiso}
$\pi_{X*}$ is affine. The relation $\equiv$ is a congruence for convex combinations, so
$\Post(X)/{\equiv}$ carries a unique convex structure making the quotient map affine, and
$\pi_{X*}$ induces an affine injection $\Post(X)/{\equiv}\hookrightarrow\D(X)$. If
$\Hyp(X)\neq\emptyset$ it is a bijection, so $\Post(X)/{\equiv}\cong\D(X)$ as convex sets.
\end{theorem}
\begin{proof}
$\pi_{X*}=\D(\mathrm{pr}_1)$ is affine. If $P_i\equiv P_i'$ then
$\pi_{X*}(\sum a_iP_i)=\sum a_i\pi_{X*}(P_i)=\pi_{X*}(\sum a_iP_i')$, so $\equiv$ is a congruence and
the induced map is well defined, affine and injective. For surjectivity pick $h_0\in\Hyp(X)$; then
$\pi_{X*}(\sum_x\rho(x)\delta_{(x,h_0)})=\rho$.
\end{proof}

\begin{proposition}[Congruence for $\oplus$]
If $P_1\equiv P_1'$ and $P_2\equiv P_2'$ then $P_1\oplus_\lambda P_2\equiv P_1'\oplus_\lambda P_2'$.
\end{proposition}

\begin{proposition}[Congruence for $\otimes$]\label{prop:cong-tensor}
If $\perp$ is total, then $\pi_{X*}(P_1\otimes P_2)=\N(m_1m_2)$ with $m_i=\pi_{X*}(P_i)$, so $\equiv$
is a congruence for $\otimes$ (no condition on $q$ is needed). If $\perp$ is not total this can fail.
\end{proposition}
\begin{proof}
With $\perp$ total, $\sum_h(P_1\cp P_2)(x,h)=\sum_{h_1,h_2}P_1(x,h_1)P_2(x,h_2)=m_1(x)m_2(x)$.
For the failure, take $X=\{x,y\}$, hypotheses $h_1,h_2,h_3$ with $h_1\perp h_3$ but not
$h_2\perp h_3$, and
$P_1=\tfrac12\delta_{(x,h_1)}+\tfrac12\delta_{(y,h_2)}$,
$P_1'=\tfrac12\delta_{(x,h_2)}+\tfrac12\delta_{(y,h_1)}$,
$P_2=\tfrac12\delta_{(x,h_3)}+\tfrac12\delta_{(y,h_3)}$. Then $P_1\equiv P_1'$ but the marginals of
$P_1\otimes P_2$ and $P_1'\otimes P_2$ are $\delta_x$ and $\delta_y$.
\end{proof}

\begin{proposition}[Congruence for $\Cond$]\label{prop:cong-cond}
Let $|X|\ge2$. Then $\equiv$ is a congruence for $\Cond_e$ (on all $P$ where both sides are defined)
if and only if $L(e\mid\cdot)$ is constant on $\Hyp(X)$.
\end{proposition}
\begin{proof}
If $L$ is constant and positive, $\Cond_e=\id$. Conversely, suppose $L(e\mid h_1)=\ell_1\ne
\ell_2=L(e\mid h_2)$, so $\ell_1+\ell_2>0$. With distinct $x_1,x_2\in X$ let
$P=\tfrac12\delta_{(x_1,h_1)}+\tfrac12\delta_{(x_2,h_2)}$ and
$P'=\tfrac12\delta_{(x_1,h_2)}+\tfrac12\delta_{(x_2,h_1)}$. Then $P\equiv P'$, but the marginals of
$\Cond_eP$ and $\Cond_eP'$ are $(\ell_1,\ell_2)/(\ell_1+\ell_2)$ and $(\ell_2,\ell_1)/(\ell_1+\ell_2)$,
which differ.
\end{proof}

\begin{remark}
In particular, requiring $L(e\mid h)=\ell(e\mid q_X(h))$ does \emph{not} make $\equiv$ a congruence for
conditioning: in the counterexample $q_X(h_1)\ne q_X(h_2)$ is allowed. With $\ell_1=2,\ell_2=1$ the
conditioned marginals are $(\tfrac23,\tfrac13)$ and $(\tfrac13,\tfrac23)$.
\end{remark}

%=====================================================================
\section{The three-level tower and the context comonad}
%=====================================================================

\subsection{The three posterior levels}

For a set $X$ let $H_3(X)=\D(X)\times\D^2(X)\times\D^3(X)$, with elements $(\rho,\Pi,\Omega)$. The
barycentric projections $b_X=\mu_X\colon\D^2X\to\D X$ and $b^{(2)}_X=\mu_{\D X}\colon\D^3X\to\D^2X$
form the tower $\D^3X\to\D^2X\to\D X$. By monad associativity
$\mu_X\circ\mu_{\D X}=\mu_X\circ\D(\mu_X)$, so the full collapse $\D^3X\to\D X$ does not depend on
which projection is used at the top. The three levels are distinct state; the projections are
observation maps.

\subsection{A general context comonad}

For a set $X$, an element $c$ of a set $A$, and $y\in X$, write $j_{X,c}(y)=(y,c)$.

\begin{theorem}[Context comonad]\label{thm:context-comonad}
Let $\E\colon\Set\to\Set$ be any functor. Put $\C_\E(X)=X\times\E(X)$ and
$\C_\E(f)=f\times\E(f)$, and define
\[
\varepsilon_X(x,c)=x,\qquad
\delta_X(x,c)=\big((x,c),\ \kappa_X(c)\big),\qquad \kappa_X(c)=\E(j_{X,c})(c)\in\E(\C_\E X),
\]
where $j_{X,c}\colon X\to\C_\E X$. Then $(\C_\E,\varepsilon,\delta)$ is a comonad on $\Set$.
\end{theorem}
\begin{proof}
Write $\C=\C_\E$. Functoriality is that of $\mathrm{id}\times\E$, and $\varepsilon$ is natural
because it is a product projection.

\emph{Naturality of $\delta$.} For $f\colon X\to Y$ we need
$\delta_Y(fx,\E f\,c)=\C\C f(\delta_X(x,c))$. The second component of the right side is
$\E(\C f)(\E(j_{X,c})(c))=\E(\C f\circ j_{X,c})(c)$. Now
$(\C f\circ j_{X,c})(y)=(fy,\E f\,c)=(j_{Y,\E f c}\circ f)(y)$, so this equals
$\E(j_{Y,\E fc})(\E f\,c)=\kappa_Y(\E f\,c)$. First components agree trivially.

\emph{Counit laws.} $\varepsilon_{\C X}\delta_X(x,c)=(x,c)$. Also
$\C(\varepsilon_X)\delta_X(x,c)=(\varepsilon_X(x,c),\E(\varepsilon_X)\E(j_{X,c})(c))=(x,c)$ since
$\varepsilon_X\circ j_{X,c}=\id_X$.

\emph{Coassociativity.} Let $w=(x,c)$ and $c'=\kappa_X(c)$. We first note
\begin{equation}\label{eq:dj}
\delta_X\circ j_{X,c}=j_{\C X,c'}\circ j_{X,c},
\end{equation}
since both send $y$ to $((y,c),c')$; this uses that $\kappa_X(c)$ depends only on $c$. Now
$\delta_{\C X}\delta_X(x,c)=\big((w,c'),\E(j_{\C X,c'})(c')\big)
=\big((w,c'),\E(j_{\C X,c'}\circ j_{X,c})(c)\big)$, while
$\C(\delta_X)\delta_X(x,c)=\big(\delta_X(w),\E(\delta_X)(c')\big)
=\big((w,c'),\E(\delta_X\circ j_{X,c})(c)\big)$. These agree by \eqref{eq:dj}.
\end{proof}

\begin{remark}
If $\E$ is constant at a set $A$, then $\C_\E$ is the reader (product) comonad $X\times A$. For general
$\E$ the context is transported into $\C_\E X$ by $\kappa_X$ rather than copied.
\end{remark}

\subsection{The three-level context}

\begin{definition}
$\E_3(X)=\D(X)\times\D^2(X)\times\D^3(X)\times\Sh(X)$ and
$\C_3=\C_{\E_3}$, so $\C_3(X)=X\times\D X\times\D^2X\times\D^3X\times\Sh X$ with points
$(x,\rho,\Pi,\Omega,\Sigma)$.
\end{definition}

\begin{corollary}\label{cor:C3}
$\C_3$ is a comonad. Explicitly, for $c=(\rho,\Pi,\Omega,\Sigma)$,
\[
\kappa_X(c)=\big((j_{X,c})_*\rho,\ \D^2(j_{X,c})(\Pi),\ \D^3(j_{X,c})(\Omega),\
\D(\Hyp(j_{X,c}))(\Sigma)\big).
\]
\end{corollary}

Each component of $\E_3(Y)$ is of the form $\D(Z)$, so $\E_3(Y)$ is a convex set, and for every $f$
the map $\E_3(f)$ is affine.

\begin{remark}[The absolute tower grows under $\D$]\label{rem:growth}
Because the levels of $\C_3$ are indexed by the object,
$\C_3(\D X)=\D X\times\D(\D X)\times\D^2(\D X)\times\D^3(\D X)\times\Sh(\D X)
=\D X\times\D^2X\times\D^3X\times\D^4X\times\Sh(\D X)$. A transformation
$\D\C_3X\to\C_3\D X$ must therefore supply a fourth-order component, and applying $\D$ again raises
the order once more. In the notation of the first version, the components $P_M\in\D X$,
$\Pi_M\in\D^2X$ and $\Omega_M\in\D^3X$ fill only the first three factors. The transport in
\cref{cor:C3} is $\rho\mapsto(j_{X,c})_*\rho$, $\Pi\mapsto\D^2(j_{X,c})\Pi$,
$\Omega\mapsto\D^3(j_{X,c})\Omega$; it is this re-indexing of the context over $\C_3X$ that makes the
order grow.
\end{remark}

\subsection{The context lift}

\begin{definition}
$J_{X,c}=\Hyp(j_{X,c})\colon\Hyp(X)\to\Hyp(\C_3X)$; explicitly
$J_{X,c}(\theta,\rho)=(\theta,(j_{X,c})_*\rho)$.
\end{definition}

\begin{proposition}[Coherence of $J$]\label{prop:CL}
\begin{enumerate}[label=(CL\arabic*)]
\item $q_{\C_3X}\circ J_{X,c}=(j_{X,c})_*\circ q_X$;
\item $\Hyp(\varepsilon_X)\circ J_{X,c}=\id_{\Hyp(X)}$;
\item $\Hyp(\delta_X)\circ J_{X,c}=J_{\C_3X,\kappa_X(c)}\circ J_{X,c}$.
\end{enumerate}
The shadow component of $\kappa_X(c)$ is $(J_{X,c})_*\Sigma$.
\end{proposition}
\begin{proof}
(CL1) is naturality of $q$. (CL2) is functoriality of $\Hyp$ and
$\varepsilon_X\circ j_{X,c}=\id$. (CL3) is functoriality of $\Hyp$ applied to \eqref{eq:dj}.
\end{proof}

%=====================================================================
\section{The distribution-level lift}
%=====================================================================

For $k\ge1$ let $\eta^{(k)}_X=\eta_{\D^kX}\circ\cdots\circ\eta_{\D X}\colon\D X\to\D^{k+1}X$, so
$\eta^{(3)}_X(\rho)=\delta_{\delta_{\delta_\rho}}$.

\begin{definition}
$K^{[k]}_X\colon\Hyp(X)\to\Hyp(\D^kX)$, $K^{[k]}_X(\theta,\rho)=(\theta,\eta^{(k)}_X(\rho))$. Then
$q_{\D^kX}(K^{[k]}_Xh)=\eta^{(k)}_X(q_Xh)$; for $k=3$ this is $\delta_{\delta_{\delta_{\rho_h}}}$.
\end{definition}

\begin{proposition}[Naturality and K/J coherence]\label{prop:K}
$K^{[k]}\colon\Hyp\Rightarrow\Hyp\circ\D^k$ is natural. In particular
\[
\Hyp(\D^3j_{X,c})\circ K^{[3]}_X=K^{[3]}_{\C_3X}\circ J_{X,c}.
\]
\end{proposition}
\begin{proof}
$\eta^{(k)}$ is a composite of whiskered units, hence natural:
$\D^{k+1}(f)\circ\eta^{(k)}_X=\eta^{(k)}_Y\circ\D(f)$. Apply with $f=j_{X,c}$.
\end{proof}

\begin{remark}
The coherence must use $\D^3j_{X,c}$: $K^{[3]}_X$ lands in $\Hyp(\D^3X)$, on which $\Hyp(\D^2j)$ does
not act. Note also that $K^{[3]}_X(h)$ is a point mass at the single point $\delta_{\delta_{\rho_h}}$ of
$\D^3X$; it is not the pushforward of $\rho_h$ along $x\mapsto\delta_{\delta_{\delta_x}}$, which would
be $\Hyp(\eta\eta\eta)$.
\end{remark}

%=====================================================================
\section{Option 1: mixed distributive laws for the absolute tower}\label{sec:laws}
%=====================================================================

Throughout, $M=\sum_ia_i\delta_{(x_i,c_i)}\in\D\C_3X$ with $c_i=(\rho_i,\Pi_i,\Omega_i,\Sigma_i)$.
The target is $\C_3\D X=\D X\times\D(\D X)\times\D^2(\D X)\times\D^3(\D X)\times\Sh(\D X)$.

\subsection{What the unit axiom forces}

\begin{lemma}\label{lem:unit-forces}
If $\lambda\colon\D\C_3\Rightarrow\C_3\D$ satisfies (U), then
$\lambda(\delta_{(x,c)})=(\delta_x,\E_3(\eta_X)(c))$. In particular its level-one component is
$\D(\eta_X)(\rho)=\sum_y\rho(y)\,\delta_{\delta_y}$, which differs from $\delta_\rho$ whenever $\rho$ is
not a point mass. The two agree after barycentric collapse: $\mu_X(\D(\eta_X)\rho)=\rho=\mu_X(\delta_\rho)$.
\end{lemma}
\begin{proof}
(U) states $\lambda(\eta_{\C_3X}(x,c))=\C_3(\eta_X)(x,c)=(\eta_Xx,\E_3(\eta_X)c)$.
\end{proof}

So a law that keeps a non-degenerate posterior $\rho$ as a single hypothesis $\delta_\rho$ at level one
cannot satisfy (U); the unit axiom requires level-one information to be spread over point
hypotheses. This is the precise form of the tension between hierarchy and barycentric collapse.

\subsection{The shift law}

\begin{definition}[Shift law]
With $K^{[1]}_X(\theta,\rho)=(\theta,\delta_\rho)$,
\[
\ls_X(M)=\Big(\sum_ia_i\delta_{x_i},\ \sum_ia_i\delta_{\rho_i},\ \sum_ia_i\delta_{\Pi_i},\
\sum_ia_i\delta_{\Omega_i},\ \sum_ia_i\,(K^{[1]}_X)_*\Sigma_i\Big).
\]
The components lie in $\D X$, $\D(\D X)$, $\D(\D^2X)=\D^2(\D X)$, $\D(\D^3X)=\D^3(\D X)$ and
$\D(\Hyp\D X)=\Sh(\D X)$ respectively, so $\ls_X\colon\D\C_3X\to\C_3\D X$.
\end{definition}

\begin{theorem}\label{thm:shift}
$\ls$ is natural and satisfies (Co). It violates (U), (M) and (Cm). The defects in (U) and (M) at
level one vanish after applying $\mu_X$.
\end{theorem}
\begin{proof}
\emph{Naturality.} Each component of $\ls$ is the pushforward of $M$ along a natural projection
(followed, for the shadow, by the natural map $K^{[1]}$ and barycenter, both natural).

\emph{(Co).} $\varepsilon_{\D X}(\ls M)=\sum_ia_i\delta_{x_i}=\D(\varepsilon_X)(M)$.

\emph{(U) fails.} By \cref{lem:unit-forces}: $\ls(\delta_{(x,c)})$ has level one $\delta_\rho$, while
$\C_3(\eta_X)(x,c)$ has $\D(\eta_X)\rho$. For $X=\{0,1\}$ and $\rho=\tfrac12\delta_0+\tfrac12\delta_1$
these are $\delta_\rho$ and $\tfrac12\delta_{\delta_0}+\tfrac12\delta_{\delta_1}$.

\emph{(M) fails.} Let $M'=\sum_\ell\nu_\ell\delta_{M_\ell}$ with $M_\ell=\sum_k a_{\ell k}
\delta_{(x_{\ell k},c_{\ell k})}$. The level-one component of $\ls\mu(M')$ is
$\sum_{\ell,k}\nu_\ell a_{\ell k}\delta_{\rho_{\ell k}}$. On the other side, $\ls_{\D X}$ reads the
level-one components $\sum_ka_{\ell k}\delta_{\rho_{\ell k}}$ of $\ls(M_\ell)$ and $\C_3(\mu_X)$ applies
$\D(\mu_X)$, giving $\sum_\ell\nu_\ell\delta_{\sum_ka_{\ell k}\rho_{\ell k}}$. These differ whenever the
$\rho_{\ell k}$ for fixed $\ell$ are not all equal, and both have barycenter
$\sum_{\ell,k}\nu_\ell a_{\ell k}\rho_{\ell k}$.

\emph{(Cm) fails.} Take $M=\tfrac12\delta_{(x,c_1)}+\tfrac12\delta_{(x,c_2)}$ with $\rho_1\neq\rho_2$.
The level-one component of $\delta_{\D X}(\ls M)$ is $\D(j_{\D X,e})(\sum_i\tfrac12\delta_{\rho_i})
=\sum_i\tfrac12\delta_{(\rho_i,e)}$, where $e$ is the pooled context of $\ls M$. On the other side,
$\C_3(\ls)\circ\ls_{\C_3X}\circ\D(\delta_X)$ gives $\sum_i\tfrac12\delta_{\ls((j_{X,c_i})_*\rho_i)}
=\sum_i\tfrac12\delta_{(\rho_i,e_i)}$ with $e_i=(\delta_{\rho_i},\delta_{\Pi_i},\delta_{\Omega_i},
(K^{[1]})_*\Sigma_i)$, and $e_i\ne e$.
\end{proof}

\begin{remark}
The transformation $\lambda^{[3]}$ of the earlier version is $\ls$ with the level-three component
omitted and a shadow component of the wrong type. Its claimed unit and comultiplication theorems
fail by \cref{thm:shift}.
\end{remark}

\subsection{The averaging law}

Let $\beta\colon\D(\E_3Y)\to\E_3Y$ be the barycenter, computed componentwise by $\mu$ on each factor
$\D(Z)$; every $\E_3(f)$ commutes with $\beta$ because it is affine.

\begin{definition}[Averaging law]
\[
\lb_X(M)=\Big(\sum_ia_i\delta_{x_i},\ \sum_ia_i\,\E_3(\eta_X)(c_i)\Big),
\]
i.e.\ $\lb_X=\big\langle\D(\varepsilon_X),\ \beta\circ\D(\E_3(\eta_X)\circ\mathrm{pr}_2)\big\rangle$.
\end{definition}

\begin{theorem}\label{thm:avg}
$\lb$ is natural and satisfies (U), (Co) and (M) exactly.
\end{theorem}
\begin{proof}
\emph{Naturality.} For $f\colon X\to Y$, using that $\E_3(\D f)$ is affine and
$\D f\circ\eta_X=\eta_Y\circ f$,
\[
\E_3(\D f)\Big(\sum_ia_i\E_3(\eta_X)c_i\Big)=\sum_ia_i\E_3(\eta_Y)(\E_3(f)c_i),
\]
which is the context component of $\lb_Y(\D(\C_3f)M)$; first components agree by naturality of
$\D(\varepsilon)$.

\emph{(U).} $\lb(\delta_{(x,c)})=(\delta_x,\E_3(\eta_X)c)=\C_3(\eta_X)(x,c)$.

\emph{(Co).} $\varepsilon_{\D X}\circ\lb_X=\D(\varepsilon_X)$ by definition.

\emph{(M).} Let $M'=\sum_\ell\nu_\ell\delta_{M_\ell}$ with
$M_\ell=\sum_ka_{\ell k}\delta_{(x_{\ell k},c_{\ell k})}$, and put
\[
P_\ell=\sum_ka_{\ell k}\delta_{x_{\ell k}},\qquad e_\ell=\sum_ka_{\ell k}\,\E_3(\eta_X)c_{\ell k}.
\]
The left side is
$\big(\sum_\ell\nu_\ell P_\ell,\ \sum_\ell\nu_\ell e_\ell\big)$. On the right,
$\D(\lb)M'=\sum_\ell\nu_\ell\delta_{(P_\ell,e_\ell)}$; then $\lb_{\D X}$ gives
\[
\Big(\sum_\ell\nu_\ell\delta_{P_\ell},\ \sum_\ell\nu_\ell\,\E_3(\eta_{\D X})e_\ell\Big),
\]
and $\C_3(\mu_X)$ turns this into
\[
\Big(\sum_\ell\nu_\ell P_\ell,\ \sum_\ell\nu_\ell\,\E_3(\mu_X\circ\eta_{\D X})e_\ell\Big)
=\Big(\sum_\ell\nu_\ell P_\ell,\ \sum_\ell\nu_\ell e_\ell\Big),
\]
using that $\E_3(\mu_X)$ is affine and $\mu_X\circ\eta_{\D X}=\id$.
\end{proof}

\begin{theorem}\label{thm:avg-comult}
$\lb$ violates (Cm).
\end{theorem}
\begin{proof}
Let $X=\{0,1\}$, fix $\Pi,\Omega,\Sigma$, and let $c_i=(\delta_{i},\Pi,\Omega,\Sigma)$ for $i=0,1$,
$M=\tfrac12\delta_{(0,c_0)}+\tfrac12\delta_{(0,c_1)}$. Write $e_i=\E_3(\eta_X)c_i$ and
$e=\tfrac12e_0+\tfrac12e_1$, so $\lb M=(\delta_0,e)$ and the level-one component of $e$ is
$\tfrac12\delta_{\delta_0}+\tfrac12\delta_{\delta_1}$. The level-one component of
$\delta_{\D X}(\lb M)$ is $\D(j_{\D X,e})$ of that, namely
$\tfrac12\delta_{(\delta_0,e)}+\tfrac12\delta_{(\delta_1,e)}$.

On the other side, $\D(\delta_X)M=\sum_i\tfrac12\delta_{((0,c_i),\kappa_X(c_i))}$. Since
$\E_3(\lb_X)$ is affine and $\lb_X\circ\eta_{\C_3X}\circ j_{X,c_i}=j_{\D X,e_i}\circ\eta_X$
(by (U)), the context component of $\C_3(\lb)\lb_{\C_3X}\D(\delta_X)(M)$ is
$\sum_i\tfrac12\E_3(j_{\D X,e_i})(e_i)$, whose level-one component is
$\tfrac12\delta_{(\delta_0,e_0)}+\tfrac12\delta_{(\delta_1,e_1)}$. Since $e_0\ne e\ne e_1$, the two sides
differ.
\end{proof}

\begin{remark}
In general the context component of the right side of (Cm) is $\sum_ia_i\E_3(j_{\D X,e_i})(e_i)$ and
that of the left side is $\E_3(j_{\D X,e})(e)$ with $e=\sum_ia_ie_i$: comultiplication copies a
\emph{single} pooled context into $\C_3\D X$, whereas the right side keeps one context per sample.
Both sides do agree after $\C_3(\varepsilon_{\D X})$ and after $\varepsilon_{\C_3\D X}$, but this
follows from (Co) and naturality alone and carries no extra information.
\end{remark}

\begin{table}[t]
\centering
\begin{tabular}{@{}lcccccc@{}}
\toprule
& natural & (U) & (Co) & (M) & (Cm) & keeps $\delta_\rho$ at level one\\
\midrule
shift law $\ls$ & yes & no & yes & no & no & yes\\
averaging law $\lb$ & yes & yes & yes & yes & no & no\\
\bottomrule
\end{tabular}
\caption{Status of the two mixed distributive laws $\D\C_3\Rightarrow\C_3\D$
(\cref{thm:shift,thm:avg,thm:avg-comult}). By \cref{lem:unit-forces}, no law satisfying (U) keeps
$\delta_\rho$ at level one.}
\label{tab:status}
\end{table}

\begin{remark}[Relation to weak distributive laws]\label{rem:weak}
The weak distributive laws of Street and B\"ohm \cite{Street2009,Bohm2010,BohmLackStreet2011}, used by
Garner \cite{Garner2020} and Goy--Petri\c{s}an \cite{GoyPetrisan2020}, are laws between two monads
that keep the multiplication axioms and weaken the unit axioms; the composite is obtained by
splitting an idempotent. Neither $\ls$ nor $\lb$ falls under that notion as stated: $\ls$ fails a
multiplication-type axiom (M), and $\lb$ fails the comonad analogue (Cm). A formulation in which a
failing axiom is witnessed by a 2-cell would need a 2-categorical structure on the base (for instance
an order or convex enrichment of the relevant hom-sets); on $\Set$ itself the only 2-cells between
natural transformations are identities, so no non-identity 2-cell $\alpha$ is available. The earlier
version's claim of a weak distributive law with a barycentric 2-cell is therefore withdrawn. An
inequality that could play the role of such a 2-cell is identified in \cref{sec:cx}.
\end{remark}

%=====================================================================
\section{Option 2: the base-relative three-level hierarchy}\label{sec:relative}
%=====================================================================

By \cref{rem:growth}, the absolute tower cannot hold exactly three levels once the monad is applied.
We now fix an outcome base $B$ (in the application, the finite outcome space of CausalIDView) and let
the hierarchy live over $B$ regardless of the object at which the comonad is evaluated.

\subsection{Coherent contexts}

\begin{definition}[Context algebra]
$A_B=\D B\times\D^2B\times\D^3B\times\Sh(B)$, with points $c=(\rho,\Pi,\Omega,\Sigma)$. It is a convex
set componentwise, and the componentwise barycenter $\beta_B\colon\D(A_B)\to A_B$ (given by $\mu$ on
each factor $\D Z$) makes $(A_B,\beta_B)$ an algebra for $\D$, namely the product of the free algebras
$(\D Z,\mu_Z)$.
\end{definition}

\begin{definition}[Coherence]
A context $c=(\rho,\Pi,\Omega,\Sigma)\in A_B$ is \emph{coherent} if
\[
\mu_B(\Pi)=\rho,\qquad \mu_{\D B}(\Omega)=\Pi,\qquad \D(q_B)(\Sigma)=\Pi .
\]
Write $A_B^{\mathrm{coh}}\subseteq A_B$ for the coherent contexts.
\end{definition}

Coherence says that the barycentric observation maps of the tower $\D^3B\to\D^2B\to\D B$ and the
predictive of the shadow agree. The top level determines the lower two, but all three remain state.

\begin{lemma}\label{lem:coh-closed}
$A_B^{\mathrm{coh}}$ is closed under $\beta_B$; hence it is a sub-algebra of $(A_B,\beta_B)$.
\end{lemma}
\begin{proof}
Each condition is an equation between affine maps ($\mu_B$, $\mu_{\D B}$, $\D(q_B)$ and the
projections), and the solution set of such an equation is closed under convex combination.
\end{proof}

\subsection{The canonical hierarchy}

\begin{definition}
$\sigma_B\colon\Sh(B)\to A_B$,
\[
\sigma_B(\Sigma)=\big(\mu_B(\D q_B\,\Sigma),\ \D q_B\,\Sigma,\ \D(\eta_{\D B})(\D q_B\,\Sigma),\ \Sigma\big).
\]
On a single hypothesis, $\sigma_B(\delta_h)=(\rho_h,\ \delta_{\rho_h},\ \delta_{\delta_{\rho_h}},\
\delta_h)$, i.e.\ $\Pi_h=\delta_{\rho_h}$ and $\Omega_h=\delta_{\Pi_h}$.
\end{definition}

\begin{proposition}\label{prop:sigma}
$\sigma_B$ is affine, injective, lands in $A_B^{\mathrm{coh}}$, and is natural in $B$:
$A_f\circ\sigma_B=\sigma_{B'}\circ\Sh(f)$ for $f\colon B\to B'$, where
$A_f=\D f\times\D^2f\times\D^3f\times\Sh f$.
\end{proposition}
\begin{proof}
Each component is a composite of affine maps, and the last component is $\Sigma$, so $\sigma_B$ is
injective. Coherence: the first condition holds by definition, the second is the unit law
$\mu_{\D B}\circ\D(\eta_{\D B})=\id$, and the third holds by definition. Naturality follows from the
naturality of $q$, $\mu$ and $\eta$.
\end{proof}

\subsection{The base-relative comonad}

\begin{definition}
For a set $Y$ put $\C^B(Y)=Y\times A_B$, $\C^B(f)=f\times\id$,
$\varepsilon_Y(y,c)=y$ and $\delta_Y(y,c)=((y,c),c)$. Write $\C^B_{\mathrm{coh}}$ for the same
construction with $A^{\mathrm{coh}}_B$.
\end{definition}

This is the product (reader) comonad with parameter $A_B$, so the comonad laws are immediate. For
$Y=B$ the underlying set $\C^B(B)=B\times\D B\times\D^2B\times\D^3B\times\Sh(B)$ coincides with
$\C_3(B)$. The difference appears under the monad:
\[
\C^B(\D Y)=\D Y\times\underbrace{\D B\times\D^2B\times\D^3B}_{\text{still three levels}}\times\Sh(B).
\]

\subsection{The strict mixed distributive law}

\begin{definition}
For $M=\sum_ia_i\delta_{(y_i,c_i)}\in\D\C^B(Y)$,
\[
\lambda^B_Y(M)=\Big(\sum_ia_i\delta_{y_i},\ \sum_ia_ic_i\Big)
=\big(\D(\mathrm{pr}_1)M,\ \beta_B(\D(\mathrm{pr}_2)M)\big)\in\C^B(\D Y),
\]
where $\sum_ia_ic_i$ is the convex combination in $A_B$, taken level by level:
\[
\textstyle\sum_ia_ic_i=\big(\sum_i a_i\rho_i,\ \sum_i a_i\Pi_i,\ \sum_i a_i\Omega_i,\ \sum_i a_i\Sigma_i\big).
\]
\end{definition}

\begin{theorem}\label{thm:relative-law}
$\lambda^B\colon\D\C^B\Rightarrow\C^B\D$ is natural and satisfies all four axioms (U), (M), (Co),
(Cm) of \cref{def:mixed}. It restricts to a mixed distributive law $\D\C^B_{\mathrm{coh}}\Rightarrow
\C^B_{\mathrm{coh}}\D$.
\end{theorem}
\begin{proof}
\emph{Naturality.} For $f\colon Y\to Y'$, $\D f\circ\D(\mathrm{pr}_1)=\D(\mathrm{pr}_1)\circ\D(f\times\id)$
and the context component is unchanged by $f\times\id$.

\emph{(U).} $\lambda^B(\delta_{(y,c)})=(\delta_y,\beta_B(\delta_c))=(\delta_y,c)=\C^B(\eta_Y)(y,c)$, by the
unit law of the algebra $\beta_B$.

\emph{(Co).} $\varepsilon_{\D Y}\circ\lambda^B_Y=\D(\mathrm{pr}_1)=\D(\varepsilon_Y)$.

\emph{(M).} First components: $\D(\mathrm{pr}_1)\circ\mu_{\C^BY}=\mu_Y\circ\D^2(\mathrm{pr}_1)$ by
naturality of $\mu$, which is the first component of $\C^B(\mu_Y)\lambda^B_{\D Y}\D(\lambda^B_Y)$.
Context components: the left side gives
$\beta_B\circ\D(\mathrm{pr}_2)\circ\mu=\beta_B\circ\mu_{A_B}\circ\D^2(\mathrm{pr}_2)$; the right side gives
$\beta_B\circ\D(\beta_B\circ\D\mathrm{pr}_2)=\beta_B\circ\D(\beta_B)\circ\D^2(\mathrm{pr}_2)$, since
$\C^B(\mu_Y)$ does not change the context. These agree by the associativity law
$\beta_B\circ\mu_{A_B}=\beta_B\circ\D(\beta_B)$ of the algebra.

\emph{(Cm).} Let $\bar c=\beta_B(\D(\mathrm{pr}_2)M)$. The left side is
$\delta_{\D Y}(\lambda^BM)=(\lambda^BM,\bar c)$. On the right,
$\D(\delta_Y)M=\sum_ia_i\delta_{((y_i,c_i),c_i)}$; applying $\lambda^B_{\C^BY}$ gives $(M,\bar c)$, and
applying $\C^B(\lambda^B_Y)$ gives $(\lambda^BM,\bar c)$.

\emph{Restriction.} $\beta_B$ preserves $A^{\mathrm{coh}}_B$ by \cref{lem:coh-closed}.
\end{proof}

\begin{remark}
\Cref{thm:relative-law} is an instance of the standard fact that a product comonad $-\times A$
distributes over a monad $T$ whenever $A$ carries a $T$-algebra structure. The content specific to
this paper is the choice of $A_B$, the coherence sub-algebra and the canonical section $\sigma_B$. By
the correspondence between mixed distributive laws and liftings \cite{PowerWatanabe2002}, $\D$ lifts
to the category of $\C^B$-coalgebras.
\end{remark}

\subsection{Pooling, hierarchy and conditioning}

\begin{proposition}[Pooling commutes with hierarchy formation]\label{prop:pooling}
If $M=\sum_ia_i\delta_{(y_i,\sigma_B(\Sigma_i))}$, then
$\lambda^B(M)=\big(\sum_ia_i\delta_{y_i},\ \sigma_B(\sum_ia_i\Sigma_i)\big)$.
\end{proposition}
\begin{proof}
$\sigma_B$ is affine (\cref{prop:sigma}).
\end{proof}

\begin{example}[Distinct posteriors survive at level two]\label{ex:level2}
Let $B=\{0,1\}$, $\rho_0=\delta_0$, $\rho_1=\tfrac13\delta_0+\tfrac23\delta_1$, and let
$M=\tfrac12\delta_{(y_0,\sigma(\delta_{h_0}))}+\tfrac12\delta_{(y_1,\sigma(\delta_{h_1}))}$ with
$q(h_k)=\rho_k$. The context of $\lambda^B(M)$ has level one
$\tfrac12\rho_0+\tfrac12\rho_1=\tfrac23\delta_0+\tfrac13\delta_1$ and level two
$\tfrac12\delta_{\rho_0}+\tfrac12\delta_{\rho_1}$: the two posteriors are kept as distinct points of
$\D B$, and level one is their barycenter, as coherence requires. In contrast, the averaging law of
the absolute tower records level one through $\D(\eta)$ (\cref{lem:unit-forces}).
\end{example}

For evidence $e$ with likelihood $L(e\mid\cdot)$ on $\Hyp(B)$, conditioning a shadow is
$\Cond_e(\Sigma)=\N(L(e\mid\cdot)\,\Sigma)$ with evidence $Z_e(\Sigma)=\sum_hL(e\mid h)\Sigma(h)$.

\begin{proposition}[Conditioning commutes with pooling]\label{prop:cond-pool}
Let $M=\sum_ia_i\delta_{(y_i,\sigma_B(\Sigma_i))}$ with $Z_i=Z_e(\Sigma_i)>0$, and let
$Z=\sum_ia_iZ_i$. Define the evidence-reweighted sample
\[
\Cond_e(M)=\sum_i\frac{a_iZ_i}{Z}\,\delta_{(y_i,\sigma_B(\Cond_e\Sigma_i))}.
\]
Then the context of $\lambda^B(\Cond_eM)$ is $\sigma_B\big(\Cond_e(\sum_ia_i\Sigma_i)\big)$.
\end{proposition}
\begin{proof}
$L\cdot\sum_ia_i\Sigma_i=\sum_ia_iZ_i\Cond_e(\Sigma_i)$, so
$\Cond_e(\sum_ia_i\Sigma_i)=\sum_i(a_iZ_i/Z)\Cond_e(\Sigma_i)$. Apply \cref{prop:pooling}.
\end{proof}

\subsection{Change of base}

\begin{proposition}\label{prop:basechange}
For $f\colon B\to B'$ the map $A_f$ is a morphism of $\D$-algebras that preserves coherence and
commutes with $\sigma$. The maps $\id_Y\times A_f$ form a comonad morphism $\C^B\Rightarrow\C^{B'}$, and
\[
(\id_{\D Y}\times A_f)\circ\lambda^B_Y=\lambda^{B'}_Y\circ\D(\id_Y\times A_f),
\]
so $\lambda^B$ is natural in the base.
\end{proposition}
\begin{proof}
$A_f$ is a product of maps $\D(g)$, hence affine, hence commutes with $\beta$. Coherence is preserved
by naturality of $\mu$ and $q$; compatibility with $\sigma$ is \cref{prop:sigma}. The comonad-morphism
equations hold because $\delta$ copies the context and $\varepsilon$ discards it. The displayed
equation follows from $A_f\circ\beta_B=\beta_{B'}\circ\D(A_f)$.
\end{proof}

\subsection{Comparison of the two options}

At $Y=B$ the two towers have the same points. They differ in the comultiplication: the absolute
comonad transports the context into distributions over $\C_3B$ through
$\rho\mapsto(j_{B,c})_*\rho$, $\Pi\mapsto\D^2(j_{B,c})\Pi$, $\Omega\mapsto\D^3(j_{B,c})\Omega$, whereas the
base-relative comonad copies it. That transport is what re-indexes the hierarchy and raises its order
under $\D$; dropping it is what allows a strict law.

\begin{table}[h]
\centering
\begin{tabular}{@{}lccccc@{}}
\toprule
& levels at $\D Y$ & (U) & (Co) & (M) & (Cm)\\
\midrule
absolute, shift law $\ls$ & four & no & yes & no & no\\
absolute, averaging law $\lb$ & four & yes & yes & yes & no\\
base-relative, $\lambda^B$ & three & yes & yes & yes & yes\\
\bottomrule
\end{tabular}
\caption{The two options. Only the base-relative tower keeps exactly three levels and admits a
strict mixed distributive law (\cref{thm:relative-law}).}
\label{tab:options}
\end{table}

%=====================================================================
\section{Coherent lifting and obstruction certificates}\label{sec:lifting}
%=====================================================================

The base-relative tower of \cref{sec:relative} supports a precise notion of when a claimed posterior
state is compatible with a hypothesis class at every level, and of what witnesses a failure. We fix a
finite base $B$ throughout.

\subsection{Admissibility data and the admissible tower}

\begin{definition}[Admissibility data]\label{def:admdata}
Admissibility data consist of a finite set $S_1\subseteq\D B$ of \emph{admissible posteriors} and a
set $S_2\subseteq\D(S_1)\subseteq\D^2B$ of \emph{admissible second-order states} such that
$\delta_\rho\in S_2$ for every $\rho\in S_1$. The \emph{admissible tower} is
\[
L_1=\operatorname{conv}(S_1)\subseteq\D B,\qquad L_2=\D(S_1)\subseteq\D^2B,\qquad
L_3=\D(S_2)\subseteq\D^3B .
\]
Here $\D(S)$ denotes the distributions supported in $S$.
\end{definition}

The basic example comes from a finite admissible hypothesis class $H\subseteq\Hyp(B)$, for instance
the hypotheses consistent with the evidence: $S_1=q_B(H)$. Two choices of $S_2$ bracket the
possibilities. In the \emph{committed regime} $S_2=\{\delta_\rho:\rho\in S_1\}$: the only admissible
second-order states are single admissible posteriors held with certainty. In the \emph{inherited
regime} $S_2=\D(S_1)$: every distribution over admissible posteriors is an admissible state. Richer
hypothesis structure, for example composite hypotheses $h_1\otimes h_2$ with an uncertain component,
supplies intermediate choices.

\begin{lemma}\label{lem:adm-tower}
Each $L_k$ is convex, $\mu_B(L_2)=L_1$ and $\mu_{\D B}(L_3)=L_2$. Moreover $\operatorname{conv}(S_2)=L_2$.
\end{lemma}
\begin{proof}
Convexity is clear. $\mu_B(\sum_\rho\Pi(\rho)\delta_\rho)=\sum_\rho\Pi(\rho)\rho$, so
$\mu_B(\D(S_1))=\operatorname{conv}(S_1)$. Since $S_2\subseteq\D(S_1)$ and $\D(S_1)$ is convex,
$\operatorname{conv}(S_2)\subseteq\D(S_1)$; conversely every $\Pi\in\D(S_1)$ is
$\sum_\rho\Pi(\rho)\delta_\rho$ with each $\delta_\rho\in S_2$. Finally
$\mu_{\D B}(\D(S_2))=\operatorname{conv}(S_2)=L_2$.
\end{proof}

\begin{proposition}[Canonical contexts realize the committed tower]\label{prop:committed}
Let $S_1=q_B(H)$ and $S_2=\{\delta_\rho:\rho\in S_1\}$. As $\Sigma$ ranges over $\D(H)$, the level
components $(\rho,\Pi,\Omega)$ of the canonical contexts $\sigma_B(\Sigma)$ range exactly over the
triples $(\mu_B\Pi,\ \Pi,\ \D(\eta_{\D B})\Pi)$ with $\Pi\in L_2$; in particular they fill $L_1$, $L_2$
and $L_3$.
\end{proposition}
\begin{proof}
$\sigma_B(\Sigma)$ has $\Pi=\D(q_B)\Sigma$, which ranges over $\D(q_B(H))=L_2$, and $\rho=\mu_B\Pi$,
$\Omega=\D(\eta_{\D B})\Pi$. The last map sends $\D(S_1)$ onto $\D(\{\delta_\rho\})=L_3$.
\end{proof}

\subsection{Claims, lifts and failure levels}

\begin{definition}[Claims and lifts]
A \emph{claim} is a partial tower $x=(x_1,x_2,x_3)$ with $x_k\in\D^kB\cup\{\ast\}$ ($\ast$ meaning
unspecified), consistent in the sense that $\mu_B(x_2)=x_1$ and $\mu_{\D B}(x_3)=x_2$ whenever both
sides are specified. A \emph{coherent lift} of $x$ is a triple $(\rho,\Pi,\Omega)\in L_1\times L_2\times
L_3$ with $\mu_B\Pi=\rho$ and $\mu_{\D B}\Omega=\Pi$ that agrees with $x$ on every specified level.
The \emph{failure level} of $x$ is
\[
\varphi(x)=\min\{k: x_k\neq\ast,\ x_k\notin L_k\}\in\{1,2,3,\infty\}.
\]
\end{definition}

\begin{theorem}[Lifting criterion]\label{thm:lifting}
Let $x$ be a claim whose highest specified level is $m$.
\begin{enumerate}[label=(\roman*)]
\item (Monotonicity) If $x_k\notin L_k$ and $x_j$ is specified for some $j>k$, then $x_j\notin L_j$.
  Hence a claim that fails at level $k$ fails at every specified level above $k$.
\item (Criterion) $x$ has a coherent lift if and only if $x_m\in L_m$, if and only if
  $\varphi(x)=\infty$.
\end{enumerate}
\end{theorem}
\begin{proof}
(i) By \cref{lem:adm-tower}, $\mu(L_{j})\subseteq L_{j-1}$, so if $x_j\in L_j$ then every specified
$x_i$ with $i<j$, being an iterated barycenter of $x_j$, lies in $L_i$. Contrapose.

(ii) If $x$ has a lift, $x_m\in L_m$. Conversely suppose $x_m\in L_m$. Below $m$ define the lift by
barycenters: $\Pi=\mu_{\D B}(x_3)$ if $m=3$, and $\rho=\mu_B(\Pi)$; these lie in $L_2$, $L_1$ by
\cref{lem:adm-tower} and agree with any specified lower $x_i$ by consistency. Above $m$ use
explicit sections: if $m\le2$ take $\Omega=\D(\eta_{\D B})(\Pi)$, which lies in $\D(\{\delta_\rho:
\rho\in S_1\})\subseteq L_3$ and satisfies $\mu_{\D B}\Omega=\Pi$; if $m=1$ first write
$x_1=\sum_\rho w_\rho\rho$ with $\rho\in S_1$ and take $\Pi=\sum_\rho w_\rho\delta_\rho\in L_2$. The
equivalence with $\varphi(x)=\infty$ follows from (i).
\end{proof}

\begin{example}[Each level can fail first]\label{ex:failures}
Let $B=\{0,1\}$, $S_1=\{\delta_0,\ \tfrac12\delta_0+\tfrac12\delta_1\}$ and the committed regime
$S_2=\{\delta_{\delta_0},\delta_{(1/2,1/2)}\}$. Then $L_1=\{(1-t,t):0\le t\le\tfrac12\}$.
\begin{enumerate}[label=(\alph*)]
\item \emph{Level one.} $x_1=\delta_1\notin L_1$.
\item \emph{Level two.} $x_2=\delta_{(3/4,1/4)}$. Its barycenter $(\tfrac34,\tfrac14)=\tfrac12\delta_0+
  \tfrac12(\tfrac12,\tfrac12)$ lies in $L_1$, but $(\tfrac34,\tfrac14)\notin S_1$, so $x_2\notin L_2$.
  The admissible $\tfrac12\delta_{\delta_0}+\tfrac12\delta_{(1/2,1/2)}$ has the same barycenter.
\item \emph{Level three.} $x_3=\delta_{\Pi'}$ with $\Pi'=\tfrac12\delta_{\delta_0}+
  \tfrac12\delta_{(1/2,1/2)}$. Then $\mu_{\D B}x_3=\Pi'\in L_2$, but $\Pi'\notin S_2$, so $x_3\notin L_3$.
  The admissible $\tfrac12\delta_{\delta_{\delta_0}}+\tfrac12\delta_{\delta_{(1/2,1/2)}}$ has the same
  barycenter.
\end{enumerate}
In each case the failure is the same move one level up: a mixture of admissible states is asserted as
a single state. At level two this is exactly the identity
$\tfrac12\delta_{\rho_1}+\tfrac12\delta_{\rho_2}\neq\delta_{(\rho_1+\rho_2)/2}$ behind the failure of the
shift law (\cref{thm:shift}).
\end{example}

\begin{proposition}[Level-three obstructions need third-order data]\label{prop:inherited}
In the inherited regime $S_2=\D(S_1)$ one has $L_3=\mu_{\D B}^{-1}(L_2)$, so no claim fails first at
level three.
\end{proposition}
\begin{proof}
$\Omega\in\D(\D(S_1))$ iff every $\Pi'\in\operatorname{supp}\Omega$ is supported in $S_1$, iff
\[
\operatorname{supp}\mu_{\D B}(\Omega)=\bigcup_{\Pi'\in\operatorname{supp}\Omega}\operatorname{supp}\Pi'
\subseteq S_1. \qedhere
\]
\end{proof}

Thus hypothesis support alone produces obstructions only at levels one and two; a genuinely
third-order obstruction requires a restriction on which second-order states are admissible.

\subsection{Obstruction certificates}

We measure obstructions by $\ell^1$ distance, $d_k(x_k)=\min_{y\in L_k}\lVert x_k-y\rVert_1$, where for
$k=2,3$ the norm is that of finitely supported measures on $\D^{k-1}B$.

\begin{theorem}[Level-one certificate]\label{thm:cert1}
$x_1\notin L_1$ if and only if there is $f\colon B\to\mathbb{R}$ with
\[
\textstyle\sum_bf(b)x_1(b)>\max_{\rho\in S_1}\sum_bf(b)\rho(b).
\]
Quantitatively,
\[
d_1(x_1)=\max_{\lVert f\rVert_\infty\le1}\Big(\langle f,x_1\rangle-\max_{\rho\in S_1}\langle f,\rho\rangle\Big),
\]
and any maximizing $f$ is an optimal certificate.
\end{theorem}
\begin{proof}
$L_1$ is a polytope in $\mathbb{R}^B$, so the first statement is the separating hyperplane theorem
(Farkas' lemma) \cite{Schrijver1986}. For the identity,
$d_1(x_1)=\min_{w\in\Delta_{S_1}}\max_{\lVert f\rVert_\infty\le1}\langle f,x_1-\sum_\rho w_\rho\rho\rangle$, since
$\max_{\lVert f\rVert_\infty\le1}\langle f,v\rangle=\lVert v\rVert_1$. The function is bilinear and both
feasible sets are polytopes, so by the minimax theorem (equivalently LP duality \cite{Schrijver1986})
the min and max may be exchanged, and $\min_w\langle f,\sum_\rho w_\rho\rho\rangle=\max_\rho\langle
f,\rho\rangle$ is attained at a vertex.
\end{proof}

\begin{proposition}[Level-two and level-three certificates]\label{prop:cert23}
$x_2\notin L_2$ iff some $\rho^\ast\in\operatorname{supp}x_2$ lies outside $S_1$, and
$d_2(x_2)=2\,x_2(\D B\setminus S_1)$. Likewise $x_3\notin L_3$ iff some
$\Pi^\ast\in\operatorname{supp}x_3$ lies outside $S_2$, and $d_3(x_3)=2\,x_3(\D^2B\setminus S_2)$.
The witnesses $\rho^\ast$ and $\Pi^\ast$ are certificates.
\end{proposition}
\begin{proof}
For $y\in\D(S)$ and outside mass $m=x(S^c)$, $\lVert x-y\rVert_1\ge\sum_{S^c}x+\big|\sum_S(x-y)\big|=
m+|(1-m)-1|=2m$. Equality holds for $y$ equal to $x$ on $S$ with the mass $m$ moved to one point of $S$.
\end{proof}

\begin{proposition}[Pooling and obstructions]\label{prop:pool-obstr}
For claims $x^{(i)}$ at the same level and weights $a_i$:
\begin{enumerate}[label=(\roman*)]
\item $d_1(\sum_ia_ix_1^{(i)})\le\sum_ia_id_1(x_1^{(i)})$, and the inequality can be strict with
  every $x_1^{(i)}\notin L_1$ and the pooled claim in $L_1$;
\item $d_k(\sum_ia_ix_k^{(i)})=\sum_ia_id_k(x_k^{(i)})$ for $k=2,3$.
\end{enumerate}
So pooling can hide level-one obstructions but never level-two or level-three ones.
\end{proposition}
\begin{proof}
(i) $d_1$ is the distance to a convex set, hence convex. For strictness take $S_1=\{(\tfrac12,\tfrac12)\}$:
$d_1(\delta_0)=d_1(\delta_1)=1$ while $\tfrac12\delta_0+\tfrac12\delta_1\in L_1$. (ii) Outside mass is
linear in the measure; apply \cref{prop:cert23}. In the same example, the level-two claim
$\tfrac12\delta_{\delta_0}+\tfrac12\delta_{\delta_1}$ has $d_2=2$.
\end{proof}

\subsection{Interpretation and scope}

\begin{definition}[Obstructed claim]
A claim is \emph{obstructed at level $k$} if $\varphi(x)=k$. Its certificate is the test function of
\cref{thm:cert1} ($k=1$) or the witness state of \cref{prop:cert23} ($k=2,3$).
\end{definition}

We propose obstructed claims as a precise candidate for what is informally called a
\emph{hallucination}: a statement that may look admissible at a lower level but admits no coherent
completion through the hierarchy generated by the admissible hypotheses. The definition is exact
and computable, since every level reduces to a linear feasibility problem with a dual certificate,
and \cref{prop:pool-obstr} suggests testing at levels two and three, where aggregation cannot mask a
failure. Whether obstruction levels predict errors in trained models is an empirical question that
this paper does not address.

\begin{remark}[Why the obstructions are convex, not homotopical]
Each $\D^kB$ is convex and hence contractible, so homotopy-theoretic obstruction classes (such as
$k$-invariants of a Postnikov tower) vanish identically on this tower, and $\mu$ has sections, so an
unconstrained lift always exists. Obstructions arise only from the admissibility constraints, and they
are convex-geometric: separating hyperplanes and support witnesses.
\end{remark}

\begin{remark}[Relation to contextuality]
When $B$ is a product of outcome sets and claims are given as marginals on several contexts, the
question whether they come from one admissible joint state is again a linear feasibility problem.
This is the setting of sheaf-theoretic contextuality \cite{AbramskyBrandenburger2011}, where failure to
admit a global section is the obstruction, and of its cohomological witnesses
\cite{AbramskyMansfieldBarbosa2012}, which are sufficient but not necessary for non-extendability. The
single-context lifting problem of this section is the base case of that framework.
\end{remark}

%=====================================================================
\section{Nondegeneracy and the likelihood-factorization criterion}
%=====================================================================

\begin{axiom}[Nondegeneracy]\label{ax:nondeg}
(a) $L(e\mid\hid)=1$ for all $e$; (b) $L(e\mid h)>0$ for all $h,e$.
\end{axiom}

\begin{remark}[Composite identifiability is too strong]\label{rem:ci}
The earlier version assumed that $\otimes$ is injective on pairs:
$h_1\otimes h_2=h_1'\otimes h_2'\Rightarrow(h_1,h_2)=(h_1',h_2')$. Together with (A2) this forces
$\Hyp(X)=\{\hid\}$: for any $h$, $h\otimes\hid=h=\hid\otimes h$, so $(h,\hid)=(\hid,h)$ and $h=\hid$.
Restricting to supports does not help when the conclusion is quantified over all posteriors.
\end{remark}

\begin{theorem}[The ``iff'']\label{thm:iff}
Assume \cref{ax:nondeg}. Then $\Cond_e(P_1\otimes P_2)=\Cond_e(P_1)\otimes\Cond_e(P_2)$ for all
$P_1,P_2$ (whenever either side is defined) if and only if
$L(e\mid h_1\otimes h_2)=L(e\mid h_1)L(e\mid h_2)$ for all compatible pairs $h_1\perp h_2$.
\end{theorem}
\begin{proof}
($\Leftarrow$) is \cref{prop:cond-comp}. ($\Rightarrow$) Let $h_1\perp h_2$. If $h_1=h_2=\hid$ there
is nothing to prove, since $L(e\mid\hid)=1$. Otherwise fix $x\in X$ and $a,b\in(0,1)$ and set
$P_1=a\delta_{(x,h_1)}+(1-a)\delta_{(x,\hid)}$, $P_2=b\delta_{(x,h_2)}+(1-b)\delta_{(x,\hid)}$. Write
$\alpha(h_1)=a$, $\alpha(\hid)=1-a$, $\beta(h_2)=b$, $\beta(\hid)=1-b$. By (A2) every pair from
$\{h_1,\hid\}\times\{h_2,\hid\}$ is compatible, and since all mass sits on one outcome, $P_1\cp P_2$
has mass $1$. For each $k\in\Hyp(X)$ put
\[
w(k)=\sum_{u\otimes v=k}\alpha(u)\beta(v)\,L(e\mid k),\qquad
w'(k)=\sum_{u\otimes v=k}\alpha(u)\beta(v)\,L(e\mid u)L(e\mid v),
\]
summing over $(u,v)\in\{h_1,\hid\}\times\{h_2,\hid\}$. Then $\Cond_e(P_1\otimes P_2)=\N(w)$ and, using
that all mass is on one outcome, $\Cond_eP_1\otimes\Cond_eP_2=\N(w')$. For pairs involving $\hid$,
$L(e\mid u\otimes\hid)=L(e\mid u)=L(e\mid u)L(e\mid\hid)$ by (A2) and \cref{ax:nondeg}(a), so
\[
w-w'=ab\,\big(L(e\mid h_1\otimes h_2)-L(e\mid h_1)L(e\mid h_2)\big)\,\delta_{h_1\otimes h_2}.
\]
The composites $h_1\otimes h_2,\ h_1,\ h_2,\ \hid$ are not all equal (else $h_1=h_2=\hid$), so there is
$k_0\neq h_1\otimes h_2$ among them, and $w(k_0)=w'(k_0)>0$ by \cref{ax:nondeg}(b). Equality
$\N(w)=\N(w')$ means $w=cw'$ for some $c>0$; evaluating at $k_0$ gives $c=1$, hence $w=w'$ and
$L(e\mid h_1\otimes h_2)=L(e\mid h_1)L(e\mid h_2)$.
\end{proof}

\begin{remark}
No identifiability or separation assumption is needed: the hypothesis $\hid$ already supplies a
reference point that fixes the normalizing constant. The earlier proof's constant
$c=Z/(Z_1Z_2)$ depended on the chosen posteriors and could not be set to $1$ by a separate choice.
\end{remark}

%=====================================================================
\section{Coalgebras, bialgebras and the 2-categorical setting}\label{sec:coalg}
%=====================================================================

This section describes the coalgebras, the lifted monad and the bialgebras of the base-relative law
$\lambda^B$ of \cref{thm:relative-law}, and makes precise which 2-category is involved.

\subsection{The ambient 2-category}\label{sec:ambient}

The monad $\D$ and the comonad $\C^B$ are a monad and a comonad on the object $\Set$ of the
2-category $\mathbf{Cat}$ of categories, functors and natural transformations, in the sense of the
formal theory of monads \cite{Street1972}. Their structure maps $\eta,\mu,\varepsilon,\delta$ and the
law $\lambda^B$ are 2-cells of $\mathbf{Cat}$, namely natural transformations between endofunctors of
$\Set$, and these 2-cells are not identities. \Cref{thm:relative-law} is therefore a statement in
$\mathbf{Cat}$, and by the correspondence of \cite{PowerWatanabe2002} it is equivalent to a lifting of
$\D$ to the category of $\C^B$-coalgebras.

The categories of coalgebras, algebras and bialgebras obtained in this way are ordinary categories.
Regarded as 2-categories (for instance by viewing each set as a discrete category), they are
\emph{locally discrete}: a natural transformation between functors of discrete categories exists only
between equal functors, so the only 2-cells are identities. \Cref{sec:2d} draws the consequences.

\subsection{Coalgebras of the base-relative comonad}

\begin{proposition}\label{prop:coalg}
A map $\gamma\colon Y\to Y\times A_B$ is a $\C^B$-coalgebra structure if and only if
$\gamma=\langle\id_Y,g\rangle$ for a (unique) map $g\colon Y\to A_B$, and a map $f\colon Y\to Y'$ is a
coalgebra morphism $(Y,\langle\id,g\rangle)\to(Y',\langle\id,g'\rangle)$ if and only if $g'\circ f=g$.
Hence $\mathrm{Coalg}(\C^B)\cong\Set/A_B$, and likewise $\mathrm{Coalg}(\C^B_{\mathrm{coh}})\cong
\Set/A^{\mathrm{coh}}_B$.
\end{proposition}
\begin{proof}
The counit law $\varepsilon_Y\circ\gamma=\id_Y$ forces the first component of $\gamma$ to be the
identity, so $\gamma=\langle\id,g\rangle$. Coassociativity then holds automatically: both
$\C^B(\gamma)\circ\gamma$ and $\delta_Y\circ\gamma$ send $y$ to $((y,g(y)),g(y))$. The condition
$(f\times\id)\circ\gamma=\gamma'\circ f$ reads $(f(y),g(y))=(f(y),g'(f(y)))$.
\end{proof}

\begin{proposition}[Cofree coalgebras]\label{prop:cofree}
The forgetful functor $U\colon\mathrm{Coalg}(\C^B)\to\Set$ has a right adjoint $F$, with
$F(X)=(X\times A_B,\delta_X)$, i.e.\ $X\times A_B$ over $A_B$ by $\mathrm{pr}_2$. Coalgebra morphisms
$h\colon(Y,\gamma)\to F(X)$ correspond to maps $\varphi\colon Y\to X$ by $h\mapsto\mathrm{pr}_1\circ h$
and $\varphi\mapsto(\varphi\times\id)\circ\gamma=\langle\varphi,g\rangle$. The counit of $U\dashv F$
at a set $X$ is $\varepsilon_X=\mathrm{pr}_1\colon X\times A_B\to X$, and the unit at a coalgebra
$(Y,\gamma)$ is the structure map $\gamma=\langle\id_Y,g\rangle\colon(Y,\gamma)\to FU(Y,\gamma)$.
\end{proposition}
\begin{proof}
Under \cref{prop:coalg}, a morphism $(Y,g)\to(X\times A_B,\mathrm{pr}_2)$ over $A_B$ is a map
$h=\langle\varphi,g\rangle$, determined by $\varphi=\mathrm{pr}_1\circ h$. The unit is the morphism
corresponding to $\id_Y$, which is $\langle\id_Y,g\rangle$.
\end{proof}

\begin{remark}
For a bare set $X$ there is no natural map $X\to X\times A_B$: a map $x\mapsto(x,e)$ requires a choice
of $e\in A_B$, and it is a coalgebra morphism only out of the constant coalgebra $(X,\mathrm{const}_e)$.
The unit of the cofree adjunction is the coalgebra structure itself, as in \cref{prop:cofree}.
\end{remark}

\subsection{The lifted monad}

\begin{proposition}\label{prop:lifted}
Under the isomorphism of \cref{prop:coalg}, the lifting of $\D$ determined by $\lambda^B$ is
\[
\tilde\D(Y,g)=\big(\D Y,\ \beta_B\circ\D(g)\big),\qquad \tilde\D(f)=\D(f),
\]
with unit and multiplication those of $\D$. Conversely, $\lambda^B_X$ is the coalgebra morphism
$\tilde\D(F X)\to F(\D X)$ corresponding under \cref{prop:cofree} to $\D(\varepsilon_X)=\D(\mathrm{pr}_1)$.
\end{proposition}
\begin{proof}
The lifting induced by a mixed law sends a coalgebra $(Y,\gamma)$ to $(\D Y,\lambda^B_Y\circ\D\gamma)$
\cite{PowerWatanabe2002}. For $\gamma=\langle\id,g\rangle$ and $P\in\D Y$,
$\lambda^B_Y(\D\gamma(P))=(P,\beta_B(\D g(P)))$, which is the coalgebra $\langle\id,\beta_B\circ\D g\rangle$.
That $\eta_Y$ and $\mu_Y$ are maps over $A_B$ is the content of (U) and (M); directly,
$\beta_B(\D g(\delta_y))=g(y)$ by the unit law of the algebra $\beta_B$, and
$\beta_B\circ\D g\circ\mu_Y=\beta_B\circ\mu_{A_B}\circ\D^2g=\beta_B\circ\D(\beta_B\circ\D g)$ by
naturality of $\mu$ and the associativity law of $\beta_B$. For the converse, $\tilde\D(FX)$ is
$\D(X\times A_B)$ over $A_B$ by $\beta_B\circ\D(\mathrm{pr}_2)$, so by \cref{prop:cofree} the morphism
corresponding to $\D(\mathrm{pr}_1)$ is $\langle\D(\mathrm{pr}_1),\beta_B\circ\D(\mathrm{pr}_2)\rangle=
\lambda^B_X$.
\end{proof}

\subsection{Bialgebras}

Recall that a $\lambda$-bialgebra is a set $Y$ with a $\D$-algebra structure $a\colon\D Y\to Y$ and a
$\C^B$-coalgebra structure $\gamma$ such that $\gamma\circ a=\C^B(a)\circ\lambda_Y\circ\D(\gamma)$.
Algebras of $\D$ are convex spaces \cite{Fritz2009}, and their morphisms are affine maps.

\begin{proposition}\label{prop:bialg}
A $\lambda^B$-bialgebra is exactly a convex space $(Y,a)$ together with an affine map
$g\colon(Y,a)\to(A_B,\beta_B)$. Hence the category of $\lambda^B$-bialgebras is isomorphic to the slice
of the category of convex spaces over $(A_B,\beta_B)$, and the coherent variant to the slice over
$(A^{\mathrm{coh}}_B,\beta_B)$.
\end{proposition}
\begin{proof}
With $\gamma=\langle\id,g\rangle$ and $P\in\D Y$, the right side of the bialgebra condition is
\[
\C^B(a)\big(P,\beta_B(\D g(P))\big)=\big(a(P),\beta_B(\D g(P))\big),
\]
and the left side is $(a(P),g(a(P)))$. So the
condition is $g\circ a=\beta_B\circ\D(g)$, which says that $g$ is a morphism of $\D$-algebras.
\end{proof}

\begin{example}
By \cref{prop:sigma}, $(\Sh(B),\mu_{\Hyp B},\sigma_B)$ is a bialgebra with values in
$A^{\mathrm{coh}}_B$: the space of shadows, mixed as distributions, together with the canonical
hierarchy it determines. For any coalgebra $(Y,g)$, the free bialgebra is
$(\D Y,\mu_Y,\beta_B\circ\D g)$. A map that reads a context off a state non-affinely, for instance from
the first hypothesis in the support of a shadow, violates the bialgebra condition.
\end{example}

So a bialgebra is a convex space of states with an affine readout of a three-level context: mixing
states mixes their contexts level by level.

\subsection{What two-dimensional monad theory does and does not add}\label{sec:2d}

Because the categories of \cref{prop:coalg,prop:lifted,prop:bialg} are locally discrete as
2-categories, every 2-cell in them, invertible or not, is an identity. Consequently, for
a 2-monad on such a 2-category, pseudo-algebras and pseudo-morphisms are strict ones, lax morphisms
are strict ones, and the inclusion of strict algebras into pseudo-algebras is an isomorphism. The
coherence and strictification results of Blackwell--Kelly--Power \cite{BKP1989,Lack2002}, their
notion of flexibility, two-dimensional monadicity \cite{Bourke2014} and doctrinal adjunction
\cite{Kelly1974} therefore reduce to statements about ordinary categories and give no information
beyond \cref{thm:relative-law} and \cref{prop:coalg,prop:cofree,prop:lifted,prop:bialg}. We make no
flexibility, strictification or 2-monadicity claims.

Genuine two-dimensional content would require three things that are not available in this paper:
\begin{enumerate}[label=(\roman*)]
\item an ambient 2-category in which contexts have non-identity morphisms, with specified
  hom-categories;
\item a 2-monad extending $\D$ to that 2-category, with a specified action on 2-cells, and a
  comonad lifted along it;
\item verification, for that 2-monad, of the codescent and doctrinal-adjunction hypotheses on which
  the results cited above depend.
\end{enumerate}

\subsection{A candidate source of 2-cells: the convex order}\label{sec:cx}

\begin{definition}[Convex order]
For $\Pi,\Pi'\in\D^2X$ with $X$ finite, write $\Pi\le_{\mathrm{cx}}\Pi'$ if
$\sum_\rho\Pi(\rho)\varphi(\rho)\le\sum_\rho\Pi'(\rho)\varphi(\rho)$ for every convex function
$\varphi\colon\D X\to\mathbb{R}$, where $\D X\subseteq\mathbb{R}^X$. Then $\mu_X\Pi=\mu_X\Pi'$, and by
Strassen's theorem \cite{Strassen1965} $\Pi'$ is obtained from $\Pi$ by a mean-preserving spread.
\end{definition}

\begin{proposition}[Shift-law defects at level one]\label{prop:cx}
For the shift law $\ls$ of \cref{thm:shift}, on level-one components:
\begin{enumerate}[label=(\alph*)]
\item (unit) $\ls(\delta_{(x,c)})$ has level one $\delta_\rho$, and
  $\delta_\rho\le_{\mathrm{cx}}\D(\eta_X)(\rho)$, the level one of $\C_3(\eta_X)(x,c)$;
\item (multiplication) with $M'$ as in the proof of \cref{thm:shift},
  $\sum_\ell\nu_\ell\,\delta_{\sum_ka_{\ell k}\rho_{\ell k}}\le_{\mathrm{cx}}
  \sum_{\ell,k}\nu_\ell a_{\ell k}\,\delta_{\rho_{\ell k}}$, i.e.\ the level one of
  $\C_3(\mu_X)\ls_{\D X}\D(\ls_X)(M')$ is below that of $\ls_X\mu_{\C_3X}(M')$.
\end{enumerate}
\end{proposition}
\begin{proof}
Both are Jensen's inequality for convex $\varphi$: in (a)
$\varphi(\rho)=\varphi(\sum_x\rho(x)\delta_x)\le\sum_x\rho(x)\varphi(\delta_x)$, and in (b)
$\varphi(\sum_ka_{\ell k}\rho_{\ell k})\le\sum_ka_{\ell k}\varphi(\rho_{\ell k})$, summed against
$\nu_\ell$.
\end{proof}

Thus at level one the unit and multiplication defects of the shift law are not arbitrary: both sides
have the same barycenter and are comparable, with pooling reducing dispersion. This suggests that
$\ls$ might be a lax mixed distributive law in a 2-category whose hom-sets are ordered by the convex
order at every level. Establishing this would require ordering all components, checking that the
functorial actions are monotone, treating the comultiplication defect, and verifying the lax
coherence axioms; we leave it open. This inequality is the precise form of the ``barycentric 2-cell''
proposed in the first version of the paper: it exists as an order relation, not as a map in $\Set$.

%=====================================================================
\section{Empirical setup}
%=====================================================================

The empirical setting is the CausalIDView benchmark of Jung, Seo, Byun, Lee and
Song~\cite{Jung2026}, which evaluates causal models across observational views; CausalIDView itself
is not gradient-free. The method evaluated here is ours: a gradient-free treatment with structured
shadows, developed in \cite[Section~51]{Awasthi-ThreeLayers}. The outcome space $X$ consists of binary causal
identification outcomes; descriptors $\theta$ are sparse multi-index hypotheses with spline links,
nominally $r=3$ directions; the posterior is sampled by gradient-free Metropolis. Targets are
Manski-type partial-identification bounds~\cite{Manski1990} and instrumental-variable bounds.

On the 40-world benchmark (\cref{app:empirical}), the posterior under a generic prior with 16 pooled
chains attains Manski RMSE $0.0934$ and IV RMSE $0.1226$, against $0.105$ and $0.138$ for the
previous structured backbone. These numbers are close to the approximate TabPFN reference
values~\cite{Hollmann2023}; since TabPFN was not rerun under the same protocol and no uncertainty
intervals are reported, we do not claim parity. The flattening of the chain-pooling curve is
consistent with, but does not by itself test, saturation of the level-two hierarchy. The 90\%
bootstrap intervals of the structured backbone cover the true bounds only 70--71\% of the time, which
indicates substantial miscalibration of that backbone. Coverage of the intervals produced by the
generic-prior posterior has not been measured; bias-aware intervals remain open.

The categorical results above do not depend on these experiments, and the experiments do not test
the categorical results.

%=====================================================================
\section{Code and data availability}
%=====================================================================

The verification scripts \texttt{check\_tower.py} (absolute tower),
\texttt{check\_base\_relative.py} (base-relative tower), \texttt{check\_lifting.py} (lifting and
obstruction certificates) and \texttt{check\_coalgebras.py} (coalgebras, bialgebras and the convex
order), and related material, are available at
\begin{center}\small\url{https://github.com/vkawasth/Trainingless-Transformer-NoGradientDescent/tree/main/Heirarchical_posteriors}\end{center}

%=====================================================================
\section{Related work}
%=====================================================================

\paragraph{Three-layer realisation, hierarchy and holonomy.}
\cite{Awasthi-ThreeLayers,Awasthi-Hierarchy,Awasthi-Holonomy}.
\paragraph{Categorical probability.} The finite distribution monad and Markov categories
\cite{Fritz2020,FGPR2023}; a synthetic de Finetti theorem \cite{FGP2021}.
\paragraph{Distributive laws.} Beck's distributive laws \cite{Beck1969}; the formal theory of monads
\cite{Street1972}; mixed laws between a monad and a comonad and their correspondence with liftings
\cite{PowerWatanabe2002}; weak distributive
laws \cite{Street2009,Bohm2010,BohmLackStreet2011}, with applications to the Vietoris monad
\cite{Garner2020} and to combining probabilistic and nondeterministic choice \cite{GoyPetrisan2020}.
\paragraph{Convex duality.} Farkas' lemma, LP duality and the minimax identity used for
obstruction certificates \cite{Schrijver1986}.
\paragraph{Information cohomology and contextuality.} Baudot--Bennequin \cite{BaudotBennequin2015},
Vigneaux \cite{Vigneaux2019}, Abramsky--Brandenburger \cite{AbramskyBrandenburger2011} and
Abramsky--Mansfield--Barbosa \cite{AbramskyMansfieldBarbosa2012}.
\paragraph{Coalgebras and two-dimensional monad theory.} Convex spaces as algebras of $\D$
\cite{Fritz2009}; two-dimensional monad theory \cite{BKP1989,Lack2002}, two-dimensional monadicity
\cite{Bourke2014} and doctrinal adjunction \cite{Kelly1974}; the convex order and Strassen's theorem
\cite{Strassen1965}.
\paragraph{Empirical.} The CausalIDView benchmark \cite{Jung2026}; the gradient-free method applied
to it \cite[Section~51]{Awasthi-ThreeLayers}; TabPFN \cite{Hollmann2023}; partial
identification \cite{Manski1990}; spiked random-matrix thresholds \cite{BBP2005}.

%=====================================================================
\section{Summary of results}
%=====================================================================

\begin{enumerate}[leftmargin=*]
\item Base-relative tower: coherent contexts form a sub-algebra; the canonical hierarchy is an
  affine, natural section; $\lambda^B$ is a strict mixed distributive law satisfying all four Beck
  axioms, natural in the base; pooling commutes with hierarchy formation and with conditioning
  (proved).
\item Coherent lifting: admissibility is convex at each level, failure is monotone, each level can
  fail first, level-three failures need third-order data, every failure has a certificate, and pooling
  hides level-one but not level-two or level-three obstructions (proved).
\item Coalgebras of the base-relative comonad are sets over $A_B$, the lifted monad is
  $(Y,g)\mapsto(\D Y,\beta_B\circ\D g)$, and bialgebras are convex spaces with an affine context map
  (proved). The ambient 2-category is $\mathbf{Cat}$; the resulting categories are locally discrete,
  so two-dimensional monad theory adds nothing here. At level one, the shift law's unit and
  multiplication defects are ordered in the convex order (proved).
\item The absolute tower grows to $\D^4X$ under $\D$ (\cref{rem:growth}).
\item $\C_\E(X)=X\times\E(X)$ is a comonad for every $\E$; $\C_3$ is an instance (proved).
\item $(\Hyp,q)$ with $q$ natural; $J=\Hyp(j)$ satisfies CL1--CL3; $K^{[k]}$ is natural and coheres
  with $J$ via $\D^3j$ (proved).
\item Normalized partial algebra on $\Post(X)$ with reweighted distributivity and conditioning laws
  (proved).
\item Affine first isomorphism theorem; congruence for $\otimes$ when $\perp$ is total; congruence for
  $\Cond_e$ iff $L(e\mid\cdot)$ is constant (proved, with counterexamples).
\item Absolute tower: shift law natural and (Co), fails (U), (M), (Cm); averaging law natural,
  (U), (Co), (M), fails (Cm). The unit axiom excludes point-mass hierarchy at level one (proved).
\item Composite identifiability trivializes $\Hyp$; the ``iff'' holds under nondegeneracy alone (proved).
\end{enumerate}

%=====================================================================
\section{Open problems}
%=====================================================================

\begin{enumerate}[leftmargin=*]
\item For the absolute tower, does any mixed distributive law $\D\C_3\Rightarrow\C_3\D$ satisfy
  all four axioms of \cref{def:mixed}? (For the base-relative tower this is settled by
  \cref{thm:relative-law}.)
\item Bayesian updating as a morphism: make $\Cond_e$ of \cref{prop:cond-pool} a natural operation
  on $\D\C^B_{\mathrm{coh}}$ compatible with $\lambda^B$, including evidence on outcomes.
\item Is there an enrichment of the base (order, convex or metric) in which the defects of $\ls$ or
  $\lb$ are witnessed by coherent 2-cells, giving a genuine weak or lax law? \Cref{prop:cx} gives the
  required inequalities for the unit and multiplication defects of $\ls$ at level one; the remaining
  components, the comultiplication defect and the lax coherence axioms are open.
\item Specify a 2-category with non-identity 2-cells and a 2-monad extending $\D$ to it, as in
  \cref{sec:2d}, and determine whether flexibility, strictification or 2-monadicity results then say
  anything about the coherent context algebra.
\item Weaken \cref{ax:nondeg} in \cref{thm:iff}, e.g.\ allow zero likelihoods.
\item Approximate admissibility: replace exact support conditions at levels two and three by
  transport distances on $\D B$ and $\D^2B$, whose Kantorovich duals would give graded certificates.
\item Empirical test: compute failure levels and certificates for model claims with known errors, and
  determine whether level-two and level-three obstructions predict them.
\item Extension to non-discrete outcome spaces (Giry monad).
\item Calibrated uncertainty and a test that could falsify the level-two saturation interpretation of
  the chain-pooling curve.
\end{enumerate}

%=====================================================================
\section{Conclusion}
%=====================================================================

Keeping exactly three posterior levels requires fixing the base on which they live. In the
absolute tower the levels are re-indexed by the object, so applying the monad raises the order to
$\D^4X$, and neither natural candidate law satisfies all Beck axioms; the unit axiom alone rules out
point-mass hierarchy at level one. In the base-relative tower the levels stay over a fixed outcome
base $B$, coherent contexts form a sub-algebra of the monad, the canonical hierarchy
$\rho_h,\ \delta_{\rho_h},\ \delta_{\delta_{\rho_h}}$ is an affine section, and pooling defines a strict
mixed distributive law satisfying all four Beck axioms. Pooling commutes with hierarchy formation and, after evidence reweighting, with conditioning,
while distinct posteriors are retained at level two. The same tower supports an exact lifting
problem: claims are checked level by level against the states a hypothesis class can produce, every
failure comes with a certificate, and obstructions above level one survive aggregation. This makes
``hallucination as obstructed lifting'' a precise, testable definition. The coalgebras and
bialgebras of the base-relative law are concrete, namely sets and convex spaces over the context
algebra; genuine two-dimensional structure is not present in this setting, and the convex order is
the most promising place to look for it.

\appendix

%=====================================================================
\section{Empirical tables}\label{app:empirical}
%=====================================================================

\subsection{The 40-world benchmark}
All numbers are on the 40-world set. The structured-backbone row uses 30 bootstrap refits per
world; the posterior rows are computed from pooled Metropolis chains (8 or 16 chains), not from
bootstrap refits. The protocols for the remaining rows are listed for confirmation in
\cref{app:changes}.

\begin{center}
\begin{tabular}{@{}lcc@{}}
\toprule
Method & Manski RMSE & IV RMSE\\
\midrule
Logistic regression (rebuild) & 0.132 & 0.169\\
Structured backbone (previous best) & 0.105 & 0.138\\
Posterior under generic prior, 8 chains & 0.0937 & 0.1230\\
Posterior under generic prior, 16 chains & 0.0934 & 0.1226\\
TabPFN, approximate reference$^{\dagger}$ & $\approx0.09$ & $\approx0.12$\\
Constant predictor & 0.131 & ---\\
\bottomrule
\end{tabular}
\end{center}
{\small $^\dagger$Approximate values, not produced under this benchmark's protocol; see
\cref{app:changes}, item~\ref{it:tabpfn}.}

\subsection{Chain-pooling curve}
\begin{center}
\begin{tabular}{@{}lccccc@{}}
\toprule
Chains & 1 & 2 & 4 & 8 & 16\\
\midrule
Manski & 0.1005 & 0.0973 & 0.0953 & 0.0937 & 0.0934\\
IV & 0.1319 & 0.1277 & 0.1250 & 0.1230 & 0.1226\\
\bottomrule
\end{tabular}
\end{center}
Differences between 8 and 16 chains (about $3\times10^{-4}$) are reported without standard errors and
should not be over-interpreted.

\subsection{Coverage}
The 90\% bootstrap intervals of the structured backbone cover the true bounds 70--71\% of the time
(82\% before the more regularized backbone). Coverage of the generic-prior posterior's intervals has
not been measured.

\subsection{Shallow-MLP worlds}
On shallow-MLP worlds the generic-prior posterior scores $0.104$ against $0.103$ for the structured
backbone, i.e.\ no improvement. A misspecified level-two shadow is one possible explanation; it has not been
tested.

\subsection{Spectral pre-pass}
The pre-pass recovers on average $2.5$ of $6.6$ planted directions and $4.5$ of $12$ active
coordinates per direction. In spiked random-matrix models, the Baik--Ben Arous--P\'ech\'e transition
\cite{BBP2005} gives a signal-strength threshold below which the top sample eigenvector carries no
information about a planted direction. Whether our spectral pre-pass operates in such a regime has
not been established, so we state the low recovery rate without attributing it to that threshold.

%=====================================================================
\section{Changes from earlier versions}\label{app:changes}
%=====================================================================

\paragraph{Version 5.} Relative to version 4 and to an intermediate draft that added two
sections on 2-categorical lifting, BKP 2-monadicity and doctrinal adjunction:
\begin{enumerate}[leftmargin=*,label=(\alph*)]
\item The intermediate sections are replaced by \cref{sec:coalg}. They worked in a 2-category of
  $\C^B$-coalgebras whose 2-cells were described as natural transformations between unspecified
  ``family actions''; no such structure had been defined, and on sets regarded as discrete categories
  the only 2-cells are identities. \Cref{sec:ambient,sec:2d} state the ambient 2-category and what it
  does and does not provide.
\item The claimed cofree coalgebra used $\Delta^B(c)=(c,e_B)$ for a chosen $e_B$, which is not the
  comultiplication $(y,c)\mapsto((y,c),c)$ of $\C^B$; the claimed adjunction unit $x\mapsto(x,e_B)$ is
  not the unit of the cofree adjunction and is not natural. Both are corrected in
  \cref{prop:cofree}. The flexibility proof, which relied on them and used an undefined map $\kappa_B$
  in the role of the barycenter $\beta_B$, is removed, as are the strictification corollary, the
  2-monadicity proposition (whose codescent and 2-cell arguments depended on the missing structure)
  and the doctrinal lax-structure formula.
\item Retained and proved: the lifted monad (\cref{prop:lifted}), with the coalgebra category
  identified as $\Set/A_B$ (\cref{prop:coalg}), and the recovery of $\lambda^B$ from the cofree
  coalgebra. New: bialgebras as convex spaces over $A_B$ (\cref{prop:bialg}) and the convex-order
  comparison of the shift law's defects (\cref{prop:cx}).
\item References: the intermediate draft cited ``Bourke and Lack, \emph{Monads and theories}, Adv.\
  Math.\ 287 (2015)''; \emph{Monads and theories} is by Bourke and Garner, and two-dimensional
  monadicity is \cite{Bourke2014}, now cited instead. A duplicate entry for \cite{Street1972} and an
  unused reference on weak bimonads are removed; broken cross-references are fixed.
\item CausalIDView is now attributed to \cite{Jung2026}; it is not gradient-free, and the
  gradient-free method applied to it is cited to \cite[Section~51]{Awasthi-ThreeLayers}.
\item Empirical corrections: the 70--71\% coverage refers to the structured backbone's bootstrap
  intervals, and coverage of the generic-prior posterior has not been measured; on shallow-MLP worlds
  $0.104$ is the generic-prior posterior and $0.103$ the structured backbone; the 30 bootstrap refits
  apply to the backbone row only, while posterior rows come from pooled Metropolis chains.
\end{enumerate}

\paragraph{Version 4.}
\begin{enumerate}[leftmargin=*,label=(\alph*)]
\item New \cref{sec:lifting}: admissibility data, admissible tower, lifting criterion, failure levels,
  obstruction certificates, pooling behavior, and the interpretation of obstructed claims.
\item New reference \cite{Schrijver1986}; abstract, introduction, summary, open problems, conclusion
  and computational checks updated.
\end{enumerate}

\paragraph{Version 3.}
\begin{enumerate}[leftmargin=*,label=(\alph*)]
\item New \cref{sec:relative}: base-relative tower with coherent contexts, canonical hierarchy
  section, and a strict mixed distributive law satisfying all four Beck axioms
  (\cref{thm:relative-law}), natural in the base, with pooling and conditioning results.
\item New \cref{rem:growth} making explicit that $\C_3(\D X)$ contains $\D^4X$, which the first
  version's $\lambda^{[3]}$ omitted; the absolute-tower laws are now labelled Option~1.
\item Abstract, introduction, summary, open problems and conclusion rewritten around the design
  choice; repository link added.
\end{enumerate}

\paragraph{Version 2 (relative to version 1).}
\begin{enumerate}[leftmargin=*]
\item Base category changed from finite sets to $\Set$, since $\D(X)$ is infinite.
\item $\Hyp$ redefined as $\Theta\times\D(-)$; transport of directions and splines by $v_j\circ f$
  removed (ill-typed, \cref{rem:variance}); the compatibility square is now naturality of $q$, not a
  definition.
\item Composition of posteriors normalized; distributivity and $\Cond$/$\otimes$ commutation
  corrected (\cref{prop:distrib,prop:cond-comp}); chain rule stated with its conditional likelihood.
\item Observational equivalence: the congruence for $\otimes$ needs total compatibility, not
  $q$-monoidality; the congruence for $\Cond$ claimed earlier is false and is replaced by
  \cref{prop:cong-cond}.
\item The comonad is now an instance of a general theorem (\cref{thm:context-comonad}); $\Sh_3$
  replaced by $\Sh$; naturality of $\delta$ is proved.
\item $J$ is defined as $\Hyp(j)$ so that CL1--CL3 are proved; the seventh tuple component of the
  earlier construction is no longer needed.
\item K/J coherence corrected to use $\D^3j$ (it was ill-typed with $\D^2j$).
\item $\lambda^{[3]}$ had four components instead of five and an ill-typed shadow; it is replaced by
  $\ls$. The earlier unit and comultiplication theorems were false (\cref{thm:shift}). The averaging
  law $\lb$ and \cref{lem:unit-forces} are new.
\item The weak distributive law and 2-cell $\alpha$ are withdrawn (\cref{rem:weak}). The contradiction
  between the claimed proof of hexagonal coherence and its listing as an open problem is resolved by
  removal.
\item Composite identifiability (which trivializes $\Hyp$) removed; it is not needed. The
  ``iff'' proof corrected (\cref{thm:iff}); condition (c) of nondegeneracy (finiteness) dropped as
  automatic.
\item Bibliography corrected. The earlier references ``B\"ohm and Lack, \emph{Weak distributive laws},
  J.~Algebra 282 (2004)'' and ``Lack, \emph{Composing monads using semimonads}, J.~Algebra 281 (2004)''
  could not be located and are removed; they are replaced by \cite{Street2009,Bohm2010,
  BohmLackStreet2011,Garner2020,GoyPetrisan2020,PowerWatanabe2002,Beck1969}. The de~Finetti reference
  now cites its journal version \cite{FGP2021}.
\item Empirical claims toned down: no parity claim with TabPFN; BBP attribution removed;
  miscalibration flagged.
\end{enumerate}

\paragraph{Items for the author to confirm.}
\begin{enumerate}[leftmargin=*,resume]
\item\label{it:tabpfn} Source and protocol of the TabPFN reference values in the benchmark table.
\item The nominal number of directions ($r=3$ in the setup) versus the $6.6$ planted directions in
  the spectral pre-pass.
\item Bootstrap standard errors for the chain-pooling curve.
\item Bibliographic details of the three self-citations
  \cite{Awasthi-Hierarchy,Awasthi-Holonomy,Awasthi-ThreeLayers}, which are reproduced as given in the
  previous version.
\item Protocols for the logistic-regression and constant-predictor rows of the benchmark table.
\item The metric (Manski or IV RMSE) behind the shallow-MLP figures $0.104$ and $0.103$.
\end{enumerate}

%=====================================================================
\section{Computational checks}
%=====================================================================

All identities and counterexamples in \cref{sec:obs,sec:laws}, \cref{thm:context-comonad} and
\cref{prop:K} were checked in exact rational arithmetic on random instances with $X=\{0,1\}$, two
structural descriptors, and maps to a three-element set (one injective, one not). The checks confirm
the comonad laws and naturality of $\varepsilon,\delta$; K/J coherence; naturality, (U), (Co) and (M)
for $\lb$; the failure of (Cm) for $\lb$ on the explicit example of \cref{thm:avg-comult}; naturality
and (Co) for $\ls$ and the failures of (U), (M), (Cm); and the counterexample of
\cref{prop:cong-cond}. These are in \texttt{check\_tower.py}.

For the base-relative tower, \texttt{check\_base\_relative.py} checks, on random coherent and
incoherent contexts over $B=\{0,1\}$: the comonad laws; naturality of $\lambda^B$ in $Y$; all four
axioms (U), (Co), (M), (Cm); closure of coherent contexts under barycenter; affineness and coherence
of $\sigma_B$; preservation of coherence and commutation with $\sigma$ and $\lambda^B$ under base change
along an injective and a non-injective map; \cref{ex:level2}; and \cref{prop:cond-pool}. All checks
pass. Both scripts are in the repository cited in the code availability section.

For \cref{sec:lifting}, \texttt{check\_lifting.py} checks the three examples of
\cref{ex:failures} with their certificates; surjectivity $\mu(L_3)=L_2$ via explicit lifts;
monotonicity of failure; the distance formula of \cref{prop:cert23} and its affinity under pooling;
the minimax identity of \cref{thm:cert1} by solving the primal and dual linear programs; the
cancellation example of \cref{prop:pool-obstr}; and that canonical contexts lie in the committed
tower. The linear programs are solved in floating point; all other checks are exact. All 17 checks
pass.

For \cref{sec:coalg}, \texttt{check\_coalgebras.py} checks that coalgebra structures are of the form
$\langle\id,g\rangle$ with coassociativity automatic; the cofree adjunction bijection; that $\eta$ and
$\mu$ are maps over $A_B$ for the lifted monad; the recovery of $\lambda^B$ from the cofree coalgebra;
the bialgebra condition for $(\Sh B,\mu,\sigma_B)$ and its failure for a non-affine context map; and
the convex-order inequalities of \cref{prop:cx} on a family of convex test functions (the proof is
Jensen's inequality). All 10 checks pass.

%=====================================================================
\begin{thebibliography}{99}
%=====================================================================

\bibitem{AbramskyBrandenburger2011}
S.~Abramsky and A.~Brandenburger.
\newblock The sheaf-theoretic structure of non-locality and contextuality.
\newblock \emph{New Journal of Physics}, 13:113036, 2011.

\bibitem{AbramskyMansfieldBarbosa2012}
S.~Abramsky, S.~Mansfield, and R.~S. Barbosa.
\newblock The cohomology of non-locality and contextuality.
\newblock \emph{Electronic Proceedings in Theoretical Computer Science}, 95:1--14, 2012.

\bibitem{Awasthi-Hierarchy}
V.~K. Awasthi.
\newblock Recovering latent hierarchy without gradient descent.
\newblock Technical report, Johns Hopkins University, September 2026.
\newblock DOI: 10.13140/RG.2.2.24981.08167.

\bibitem{Awasthi-Holonomy}
V.~K. Awasthi.
\newblock Holonomy obstructions for composed local operations.
\newblock Technical report, Johns Hopkins University, September 2026.
\newblock DOI: 10.13140/RG.2.2.25336.20483.

\bibitem{Awasthi-ThreeLayers}
V.~K. Awasthi.
\newblock Three layers of realisation: hierarchy recovery, holonomy, information and contextuality.
\newblock Technical report, Johns Hopkins University, September 2026.
\newblock DOI: 10.13140/RG.2.2.24274.52163.

\bibitem{BKP1989}
R.~Blackwell, G.~M. Kelly, and A.~J. Power.
\newblock Two-dimensional monad theory.
\newblock \emph{Journal of Pure and Applied Algebra}, 59(1):1--41, 1989.

\bibitem{Bourke2014}
J.~Bourke.
\newblock Two-dimensional monadicity.
\newblock \emph{Advances in Mathematics}, 252:708--747, 2014.

\bibitem{BBP2005}
J.~Baik, G.~Ben~Arous, and S.~P\'ech\'e.
\newblock Phase transition of the largest eigenvalue for nonnull complex sample covariance matrices.
\newblock \emph{Annals of Probability}, 33(5):1643--1697, 2005.

\bibitem{BaudotBennequin2015}
P.~Baudot and D.~Bennequin.
\newblock The homological nature of entropy.
\newblock \emph{Entropy}, 17(5):3253--3318, 2015.

\bibitem{Beck1969}
J.~Beck.
\newblock Distributive laws.
\newblock In \emph{Seminar on Triples and Categorical Homology Theory}, volume~80 of \emph{Lecture
  Notes in Mathematics}, pages 119--140. Springer, 1969.

\bibitem{Bohm2010}
G.~B\"ohm.
\newblock The weak theory of monads.
\newblock \emph{Advances in Mathematics}, 225(1):1--32, 2010.
\newblock doi:10.1016/j.aim.2010.02.015.

\bibitem{BohmLackStreet2011}
G.~B\"ohm, S.~Lack, and R.~Street.
\newblock On the 2-categories of weak distributive laws.
\newblock \emph{Communications in Algebra}, 39(12):4567--4583, 2011.
\newblock doi:10.1080/00927872.2011.616436.

\bibitem{Fritz2020}
T.~Fritz.
\newblock A synthetic approach to {M}arkov kernels, conditional independence and theorems on
  sufficient statistics.
\newblock \emph{Advances in Mathematics}, 370:107239, 2020.

\bibitem{Fritz2009}
T.~Fritz.
\newblock Convex spaces {I}: definition and examples.
\newblock arXiv:0903.5522, 2009.

\bibitem{FGP2021}
T.~Fritz, T.~Gonda, and P.~Perrone.
\newblock De {F}inetti's theorem in categorical probability.
\newblock \emph{Journal of Stochastic Analysis}, 2(4):Article~6, 2021.
\newblock doi:10.31390/josa.2.4.06. arXiv:2105.02639.

\bibitem{FGPR2023}
T.~Fritz, T.~Gonda, P.~Perrone, and E.~F. Rischel.
\newblock Representable {M}arkov categories and comparison of statistical experiments in categorical
  probability.
\newblock \emph{Theoretical Computer Science}, 961:113896, 2023.

\bibitem{Garner2020}
R.~Garner.
\newblock The {V}ietoris monad and weak distributive laws.
\newblock \emph{Applied Categorical Structures}, 28(2):339--354, 2020.
\newblock doi:10.1007/s10485-019-09582-w.

\bibitem{GoyPetrisan2020}
A.~Goy and D.~Petri\c{s}an.
\newblock Combining probabilistic and non-deterministic choice via weak distributive laws.
\newblock In \emph{Proceedings of the 35th Annual ACM/IEEE Symposium on Logic in Computer Science
  (LICS)}, pages 454--464, 2020.

\bibitem{Hollmann2023}
N.~Hollmann, S.~M\"uller, K.~Eggensperger, and F.~Hutter.
\newblock {TabPFN}: A transformer that solves small tabular classification problems in a second.
\newblock In \emph{International Conference on Learning Representations (ICLR)}, 2023.

\bibitem{Jung2026}
H.~Jung, G.~Seo, H.~Byun, J.~Lee, and K.~Song.
\newblock What you observe determines how you identify causal effects: evaluating causal models
  across observational views.
\newblock arXiv:2609.36881, 2026.

\bibitem{Kelly1974}
G.~M. Kelly.
\newblock Doctrinal adjunction.
\newblock In \emph{Category Seminar (Sydney, 1972/1973)}, volume 420 of \emph{Lecture Notes in
  Mathematics}, pages 257--280. Springer, 1974.

\bibitem{Lack2002}
S.~Lack.
\newblock Codescent objects and coherence.
\newblock \emph{Journal of Pure and Applied Algebra}, 175(1--3):223--241, 2002.

\bibitem{Manski1990}
C.~F. Manski.
\newblock Nonparametric bounds on treatment effects.
\newblock \emph{American Economic Review}, 80(2):319--323, 1990.

\bibitem{PowerWatanabe2002}
J.~Power and H.~Watanabe.
\newblock Combining a monad and a comonad.
\newblock \emph{Theoretical Computer Science}, 280(1--2):137--162, 2002.
\newblock doi:10.1016/S0304-3975(01)00024-X.

\bibitem{Schrijver1986}
A.~Schrijver.
\newblock \emph{Theory of Linear and Integer Programming}.
\newblock Wiley, Chichester, 1986.

\bibitem{Strassen1965}
V.~Strassen.
\newblock The existence of probability measures with given marginals.
\newblock \emph{Annals of Mathematical Statistics}, 36(2):423--439, 1965.

\bibitem{Street1972}
R.~Street.
\newblock The formal theory of monads.
\newblock \emph{Journal of Pure and Applied Algebra}, 2(2):149--168, 1972.

\bibitem{Street2009}
R.~Street.
\newblock Weak distributive laws.
\newblock \emph{Theory and Applications of Categories}, 22(12):313--320, 2009.

\bibitem{Vigneaux2019}
J.~P. Vigneaux.
\newblock \emph{Topology of Statistical Systems: A Cohomological Approach to Information Theory}.
\newblock PhD thesis, Universit\'e Paris Diderot, 2019.

\end{thebibliography}

\end{document}
